\part{The exact inverse in general dimension}\label{gd:part}
% Certified data are loaded once from gd_data_block.tex in the preamble.


\section{The general-dimensional certificate}\label{gd:sec:theorem}

This part proves the cases $n\in\gdDimensions\setminus\{5\}$ of
Theorem~\ref{thm:main} and gives a second proof of the case $n=5$.

\begin{theorem}\label{gd:thm:main}
For every $n\in\gdDimensions$ there is a bounded strictly convex domain
$\Omega\subset\R^n$, different from a ball and strictly star-shaped
with respect to the origin, with real-analytic boundary
diffeomorphic to $S^{n-1}$, invariant under $O(n-1)$ on the last $n-1$
coordinates and under reflection of the first coordinate. It carries a
nonconstant real-valued solution, real analytic on a neighbourhood of
$\overline\Omega$, of
\begin{equation}\label{gd:eq:schiffer}
 (\Delta+1)u=0\quad\hbox{in }\Omega,\qquad
 u=1,\quad \partial_{\mathbf n} u=0\quad\hbox{on }\partial\Omega.
\end{equation}
Here $\mathbf n$ is the outward unit normal.
In particular, $\Omega$ fails the Pompeiu property.
\end{theorem}

For $n=5$ this is a second proof of the five-dimensional case of
Theorem~\ref{thm:main}; it uses none of the bounds of Section~\ref{r5:sec}.
The argument below is parametrised by an integer $n\ge5$.
The angular conversion uses $\lambda\ge1$, and root coverage uses
$\nu=l+\lambda\ge3/2$.
Its exact parameters and constants are in
Tables~\ref{gd:tab:inputs}--\ref{gd:tab:geometry}; the finite centres are
the compressed files \path{anc/gd/frozen_centers/nN.json.gz} described in
Subsection~\ref{gd:sec:verification}.
The scale $\rho$, analytic weight $R>1$, shape weight $\sigma>0$ and
existence radius $r_*>0$ are exact, as are $\lambda$, $L_0$ and the mode
and state counts. Entries carrying an inequality sign are strict
one-sided bounds.

\begin{table}[htbp]
\centering\small
\def\gdRow#1#2#3#4#5#6#7#8{$#1$ & #2\\}
\begin{tabular}{@{}crrrrr@{}}
\toprule
$n$ & stored modes & $\rho$ & $R$ & $\sigma$ & $r_*$\\
\midrule
\gdData
\bottomrule
\end{tabular}
\caption{Exact inputs. The stored-mode count is the length of each finite
shape and field array, not a count of converged coefficients; the clamping
lift and the solution have infinite expansions.}\label{gd:tab:inputs}
\end{table}

\subsection{Meridian coordinates and the equation}\label{gd:sec:formulation}
Write reference points as $x=(x_1,x')\in\R\times\R^{n-1}$ and set
\[
 y=|x'|,\quad w=x_1+iy,\quad r=|x|,\quad \xi=x_1/r,\quad
 a=n-2,\quad\lambda=a/2,\quad\mu=\rho^2.
\]
As in Sections~\ref{r5:sec} and~\ref{r7:sec}, $\xi$ is the angular
variable. The radial variable here is $r$; below $t$ denotes a time or
a parameter.
All scalar unknowns are $O(n-1)$-invariant and even in $x_1$.
The regular lift to the full ball interprets the meridian operator
\begin{equation}\label{gd:eq:laplacian}
 \Delta_n=\partial_{x_1}^2+\partial_y^2+\frac a y\partial_y .
\end{equation}
There is no boundary condition on the axis.

Let $\psi$ be odd holomorphic with real coefficients, $p=\psi'$ and
$q=\im\psi/y$. As in Lemma~\ref{r5:lem-lift}, put
\begin{equation}\label{gd:eq:lift-map}
 \Theta_\psi(x_1,x')=
 \bigl(\re\psi(x_1+i|x'|),\,q(x_1,|x'|)x'\bigr).
\end{equation}
Hereafter $b$ is the positive coefficient of $w$ in the fixed centre
$\psi_0$; it is held constant throughout the equation.
The unscaled domain is $\Theta_\psi(B^n)$, with $u_0$ of frequency
$\rho$; the domain of the theorem is $\Omega=\rho\Theta_\psi(B^n)$.
On the meridian plane the differential is $|p|$ times a rotation;
on each of the $n-2$ transverse directions it is multiplication by $q$.
Consequently
\[
 \det D\Theta_\psi=|p|^2q^{n-2}.
\]
The proof of Lemma~\ref{r5:lem-lift}, with $n-2$ transverse vectors,
gives the change of variables in every dimension.

For $U=u_0\circ\Theta_\psi$, with frequency $\rho$, the
pull-back equation is
\[
 \operatorname{div}(q^a\nabla U)+\mu q^a|p|^2U=0.
\]
Divide by the positive factor $bq^{a-1}$ and define
\begin{equation}\label{gd:eq:B}
 \begin{split}
 BU&=\frac1b\left\{q\Delta_n U+
 a\left(q_{x_1}U_{x_1}+\frac{q_y}{y}\,yU_y\right)
 +\mu q|p|^2U\right\},\\
 \mathcal L U&=\Delta_n U+\frac a q\nabla q\cdot\nabla U
                  +\mu|p|^2U,\qquad B=(q/b)\mathcal L .
 \end{split}
\end{equation}
In particular the multiplier relating $B$ to the operator $\Delta+\mu$ is
\begin{equation}\label{gd:eq:physical-multiplier}
 B(u_0\circ\Theta_\psi)=
 \frac{q|p|^2}{b}\bigl((\Delta+\mu)u_0\bigr)\circ\Theta_\psi .
\end{equation}
The quotients in~\eqref{gd:eq:B} are regular analytic scalars:
$q=\int_0^1\re p(x_1+isy)\,ds$ is even in $y$, and so is $q_y/y$.
This is the quartic normalisation of~\eqref{r5:eq-div}, with the
fixed scalar factor $1/b$.

\section{Angular coefficient spaces and the compatible Poisson inverse}
\label{gd:sec:spaces}

\subsection{The angular algebra}\label{gd:sec:algebra}
Set
\[
 G_l(\xi)=\frac{C_l^\lambda(\xi)}{C_l^\lambda(1)},\qquad
 H_l(x)=r^lG_l(\xi).
\]
For $n=5,7$ these are the polynomials denoted $Y_\ell,H_\ell$ in
Sections~\ref{r5:sec-spaces} and~\ref{r7:sec-spaces}.
The generating function of the Gegenbauer polynomials gives
\begin{equation}\label{gd:eq:cosine}
 \begin{split}
 G_l(\cos\theta)&=\sum_mD_{ml}\cos(m\theta),\\
 D_{l-2k,l}&=
 \frac{(2-\delta_{l,2k})(\lambda)_k(\lambda)_{l-k}}
 {k!(l-k)!C_l^\lambda(1)}
 \quad(0\le k\le l/2),\qquad
 h_l=\sum_m D_{ml}R^m .
 \end{split}
\end{equation}
In particular $D_{ml}\ge0$, $\sum_mD_{ml}=1$ and
$h_l\ge\sum_mD_{ml}=1$.

Let $\mathcal Y_n$ consist of expansions
$f=\sum_{l\ge0}f_l(r)G_l(\xi)$, with $f_l\in C([0,1])$,
$f_l(0)=0$ for $l>0$, and norm
\begin{equation}\label{gd:eq:Y}
 \nrm f_{\mathcal Y_n}=\sum_{l\ge0}h_l\nrm{f_l}_\infty.
\end{equation}
Let $Y$ be its real even sector, so only even $l$ occur.
The radial norm is a supremum over all radii. It replaces the sum of
Jacobi radial coefficients, graded by polynomial degree, in
Sections~\ref{r5:sec-spaces} and~\ref{r7:sec-spaces}.
In particular a continuous radial profile is admitted without a Jacobi
expansion. The Poisson estimates below provide its $C^1$ reconstruction.

The raw cosine norm is
$\nrm f_{\rm raw}=\sum_mR^m\nrm{\tilde f_m}_\infty$ for the corresponding
cosine expansion. The boundary norms $W_0,W_1$ have exactly the Fourier
convention of~\eqref{r2:eq:Wone}, now with weight $R^{|m|}$.
For a real cosine series they are
\[
 \nrm h_{W_0}=\sum_{m\ge0}R^m|h_m|,\qquad
 \nrm h_{W_1}=\sum_{m\ge0}(1+m)R^m|h_m|.
\]
Thus $\nrm f_{\rm raw}\le\nrm f_{\mathcal Y_n}$ and boundary evaluation
has norm at most one into $W_0$.

\begin{lemma}\label{gd:lem:algebra}
The spaces $\mathcal Y_n$ and $Y$ are Banach algebras with
$\nrm{fg}_{\mathcal Y_n}\le\nrm f_{\mathcal Y_n}\nrm g_{\mathcal Y_n}$.
\end{lemma}
\begin{proof}
Completeness is the weighted $C([0,1])$ argument of
Lemma~\ref{r2:lem:algebra}; the origin conditions define a closed subspace.
Since $|G_l|\le1$, convergence in the norm gives uniform convergence on
the full closed ball. For the product, write
$G_lG_k=\sum_dL_{lk}^dG_d$. The addition theorem represents $G_l$
as a positive semidefinite zonal kernel on $S^{n-1}$.
The product kernel is positive semidefinite by tensoring its feature
vectors. Its scalar eigenvalues on spherical harmonics are nonnegative,
and orthogonality identifies them, up to positive normalisations, with
$L_{lk}^d$. Hence $L_{lk}^d\ge0$; evaluation at $\xi=1$ gives
$\sum_dL_{lk}^d=1$. This is the positivity used in
Theorem~\ref{r5:thm-mult}(a), also given by the Gegenbauer linearisation
formula~\cite[18.18.22]{DLMF}.
The nonnegative cosine expansions and
$T_mT_j=(T_{m+j}+T_{|m-j|})/2$ give
\[
 \sum_dL_{lk}^dh_d\le h_lh_k .
\]
Taking radial suprema proves the product estimate for finite sums.
Absolute convergence proves it for infinite sums. At the origin only
constant angular terms remain, so products preserve origin compatibility.
\end{proof}

\begin{lemma}[Conversion of the full cosine expansion]\label{gd:lem:connection}
Put $z=R^{-2}$ and $e_i=|(-\lambda)_i|/i!$. Then
\begin{equation}\label{gd:eq:kappa}
 \nrm f_{\mathcal Y_n}\le\kappa\nrm f_{\rm raw},\qquad
 \kappa=(1-z)^{-\lambda}\sum_{i\ge0}e_i z^i .
\end{equation}
This bound applies to arbitrary continuous radial profiles and includes
all angular degrees.
\end{lemma}
\begin{proof}
The connection formula of Lemma~\ref{r7:lem-conn}, whose derivation of
the coefficient formula holds for every $\lambda>0$, is
\begin{equation}\label{gd:eq:connection}
 T_N=\sum_{i=0}^{\lfloor N/2\rfloor} C_{N-2i,N}G_{N-2i},\qquad
 C_{l,N}=
 \frac{N(-\lambda)_i(N-i-1)!(l+\lambda)C_l^\lambda(1)}
 {2i!(\lambda)_{N-i+1}},
 \quad l=N-2i,
\end{equation}
for $N>0$; $T_0=G_0=1$. It follows by the same limit of the
Gegenbauer connection formula~\cite[18.18.16,18.7.25]{DLMF}.
For $l>0$ the top cosine coefficient is
$2(\lambda)_l/(l!C_l^\lambda(1))$.
Since $\lambda\ge1$, the remaining positive coefficients in
\eqref{gd:eq:cosine} give
\[
 h_l\le
 \frac{2(\lambda)_lR^l}{l!C_l^\lambda(1)}(1-z)^{-\lambda}.
\]
For $l=0$ use the top cosine coefficient one and
$h_0\le(1-z)^{-\lambda}$.
For $i\ge1$, multiplication by $|C_{l,N}|$ leaves
$e_iR^Nz^i(1-z)^{-\lambda}$ times
\[
 \frac{N}{N-i+\lambda}
 \prod_{s=1}^{i-1}\frac{l+s}{l+s+\lambda}\le1.
\]
For $i=0$, the product of $C_{N,N}$ and the top cosine
coefficient is exactly one. For $i\ge1$ the displayed product is at most
$(l+1)/(l+i)$, and the desired inequality reduces to
$i^2-\lambda i+N(\lambda-1)\ge i(i+\lambda-2)\ge0$,
using $N\ge2i$. Summation proves~\eqref{gd:eq:kappa}.
Conversion lowers degree and preserves the origin condition.

For integer $\lambda$, the infinite series is $(1+z)^\lambda$.
For $\lambda=m+1/2$, the Taylor coefficients $(-\lambda)_i/i!$
have sign $(-1)^i$ for $i\le m+1$ and fixed sign
$\epsilon_\lambda=(-1)^{m+1}$ for $i\ge m+1$. Hence
\[
 \sum_{i\ge0}e_i z^i=\epsilon_\lambda(1-z)^\lambda
 +2\sum_{\substack{0\le i\le m\\i\equiv m\ (\mathrm{mod}\ 2)}}e_i z^i.
\]
For $\lambda=3/2,5/2$ this is the expression in
Lemmas~\ref{r5:lem-conn}(b) and~\ref{r7:lem-conn}.
Thus~\eqref{gd:eq:kappa} has a finite closed expression in both cases.
\end{proof}

\begin{lemma}[Weights and polynomial moments]\label{gd:lem:moments}
One has $1=h_0\le h_1\le h_2\le\cdots$ and $h_l\le R^l$.
For a representation $f=\sum_{l,s}c_{ls}r^{2s}H_l$ define
$\mathfrak m_k(f)=\sum_{l,s}|c_{ls}|h_l(l+2s)^k$.
Then
\begin{equation}\label{gd:eq:poly-derivatives}
 \nrm{\partial_{x_1}f}_{\mathcal Y_n}\le C_x^{\rm pol}\mathfrak m_1(f),
 \qquad
 \nrm{y\partial_yf}_{\mathcal Y_n}\le C_y^{\rm pol}\mathfrak m_1(f),
 \quad C_x^{\rm pol}=1+R,\quad C_y^{\rm pol}=1+RC_x^{\rm pol}.
\end{equation}
\end{lemma}
\begin{proof}
Choose a random $\zeta$ with distribution $\operatorname{Beta}(\lambda,\lambda)$,
then independent signs of mean $2\zeta-1$, and let $S_l$ be their sum.
Integration of the binomial law gives
$\binom lk(\lambda)_k(\lambda)_{l-k}/(2\lambda)_l$.
Pairing $k$ and $l-k$ identifies~\eqref{gd:eq:cosine} with
$h_l=\mathbb E R^{|S_l|}$.
The conditional expectation is
$\mathbb E(S_{l+1}\mid S_l)=S_l(1+(2\lambda+l)^{-1})$.
Jensen's inequality for the even convex function $s\mapsto R^{|s|}$
proves monotonicity; $|S_l|\le l$ proves the upper bound.
The harmonic identity $\partial_{x_1}H_l=lH_{l-1}$ is proved in
Section~\ref{r7:sec-spaces} by an argument valid in every dimension.
Differentiating $r^{2s}H_l$ and using
$\nrm{x_1}_{\mathcal Y_n}=R$ gives the first estimate.
The Euler derivative multiplies that monomial by $l+2s$;
$y\partial_y=r\partial_r-x_1\partial_{x_1}$ gives the second.
Absolute convergence with the indicated moments justifies all derivatives.
\end{proof}

\subsection{The zero-Dirichlet inverse and its trace}\label{gd:sec:Poisson}
Let $K_n^D$ be the regular zero-Dirichlet inverse of $\Delta_n$.
For angular degree $l$ its radial equation is
\[
 \mathcal D_lv=v''+\frac{n-1}{r}v'-\frac{l(l+n-2)}{r^2}v=f_l,
 \qquad v(1)=0.
\]
Regularity at the origin is that of the solution on $B^n$.

\begin{lemma}\label{gd:lem:Poisson}
For $f\in Y$ the following bounds hold:
\begin{equation}\label{gd:eq:Poisson-constants}
 \begin{split}
 \nrm{K_n^Df}_Y&\le C_K\nrm f_Y,\qquad C_K=\frac1{2n},\\
 \nrm{\partial_{x_1}K_n^Df}_{\mathcal Y_n}&\le C_x^K\nrm f_Y,\qquad
 \nrm{y\partial_yK_n^Df}_Y\le C_y^K\nrm f_Y,\\
 C_x^K&=\frac Rn+\frac R{n+2}+\frac{2(2R+1)}{n+1},\\
 C_E^K&=\frac1n+\frac1{n+2}+\frac2{n+1},
 \qquad C_y^K=C_E^K+RC_x^K.
 \end{split}
\end{equation}
Moreover,
\begin{equation}\label{gd:eq:Poisson-trace}
 (K_n^Df)_l'(1)=\int_0^1r^{l+n-1}f_l(r)\,dr,\qquad
 \nrm{\partial_rK_n^Df}_{W_1}\le\nrm f_Y.
\end{equation}
\end{lemma}
\begin{proof}
For $-\mathcal D_l$ the function $(1-r^2)/(2n)$ is a supersolution
with right side at least one. The radial maximum principle therefore gives
$\nrm v_\infty\le\nrm{f_l}_\infty/(2n)$.
For $l\ge2$, $-\mathcal D_l r=(l-1)(l+n-1)/r$; comparison with
$r\nrm{f_l}_\infty/d_l$, where $d_l=(l-1)(l+n-1)$, gives
$|v(r)|/r\le\nrm{f_l}_\infty/d_l$.
For $l=0$, apply the maximum principle on $B^n$ to the radial functions
$(1-r^2)\nrm{f_0}_\infty/(2n)\pm v$. For $l\ge1$, apply it on
$[\epsilon,1]$ and let $\epsilon\to0$. Continuity gives $v(\epsilon)=o(1)$ for $l\ge1$.
Full-ball Dirichlet regularity for continuous data~\cite{GT01} gives a $C^1$
solution; its constant and linear Taylor terms have angular degrees
zero and one. Thus $v(\epsilon)=o(\epsilon)$ for $l\ge2$, as
required for the second barrier.
Writing $v=r^l\phi$ and integrating its equation gives
\begin{equation}\label{gd:eq:radial-derivative}
 v'=\frac{lv}{r}+\mathcal I_l(r),\qquad
 \mathcal I_l(r)=r^{-l-n+1}\int_0^r s^{l+n-1}f_l(s)\,ds,\qquad
 |\mathcal I_l(r)|\le\frac r{l+n}\nrm{f_l}_\infty .
\end{equation}
For $l=0$ this gives $|v'|\le\nrm{f_0}_\infty/n$.

The Gegenbauer identity
$(1-\xi^2)G_l'=l(G_{l-1}-\xi G_l)$, obtained by differentiating
$H_l$, gives the cancellation
\[
 \partial_{x_1}(vG_l)=\xi\mathcal I_lG_l+(lv/r)G_{l-1}.
\]
Multiplication by $\xi=G_1$ costs $R$.
For even $l\ge2$, $l/d_l\le2/(n+1)$, and $h_{l-1}\le h_l$.
Thus the coefficient bound for $\partial_{x_1}K_n^D$ is at most
$\max\{R/n,R/(n+2)+2/(n+1)\}$.
For the Euler derivative $r\partial_r$ the corresponding bound is
$\max\{1/n,1/(n+2)+2/(n+1)\}$.
The enlarged constants in~\eqref{gd:eq:Poisson-constants} dominate
these bounds, and $y\partial_y=r\partial_r-x_1\partial_{x_1}$
gives $C_y^K$. Inputs have even degree, so no $l=1$ radial estimate
is required; $\partial_{x_1}$ may have odd outputs in $\mathcal Y_n$.

At $r=1$, \eqref{gd:eq:radial-derivative} and $v(1)=0$
give the moment formula. The boundary cosine frequencies of $G_l$
are at most $l$, hence
\[
 \nrm{\partial_rK_n^Df}_{W_1}
 \le\sum_l(l+1)h_l|v_l'(1)|
 \le\sum_l\frac{l+1}{l+n}h_l\nrm{f_l}_\infty\le\nrm f_Y .
\]
For continuous data the full-ball Poisson solution belongs to
$W^{2,p}(B^n)$ for every finite $p$ and to $C^1(\overline{B^n})$
for $p>n$, by Dirichlet regularity~\cite{GT01}. Approximation of the data proves the moment formula
for its classical normal trace.
\end{proof}

Define the closed compatible subspace and product space
\begin{equation}\label{gd:eq:X}
 \begin{split}
 \Gc&=\left\{g\in Y:\int_0^1r^{l+n-1}g_l(r)\,dr=0
                            \text{ for every even }l\right\},\\
 K_n&=K_n^D|_{\Gc},\qquad
 W=\left\{p=\sum_{j\ge0}p_{2j}w^{2j}:
                  p_{2j}\in\R,\ \nrm p_W=\sum_{j\ge0}R^{2j}|p_{2j}|<\infty\right\},\\
 X&=\Gc\times W,\qquad
 \nrm{(g,p)}_X=\nrm g_Y+\sigma\nrm p_W .
 \end{split}
\end{equation}
The moments are bounded functionals, so $\Gc$ is complete.
Lemma~\ref{gd:lem:Poisson} gives both zero traces of $K_ng$.
The normalised primitive is $\psi(w)=\int_0^wp(z)\,dz$.
The nonlinear equation is
\begin{equation}\label{gd:eq:F}
 F(g,p)=B(1+K_ng)=0 .
\end{equation}
It is affine in $g$ and cubic in $p$, thus a polynomial of degree four
in $(g,p)$.
The constant $\sigma$ changes the product norm, not this equation.

\begin{table}[htbp]
\centering\small
\def\gdRow#1#2#3#4#5#6#7#8{$#1$ & #3\\}
\begin{tabular}{@{}crrrrr@{}}
\toprule
$n$ & $\lambda$ & $\kappa$ & $C_x^K$ & $C_y^K$ & $L_0$\\
\midrule
\gdData
\bottomrule
\end{tabular}
\caption{Angular conversion, Poisson bounds, and the starting degree
$L_0$ of the high angular complement. Here $C_K=1/(2n)$ exactly.}
\label{gd:tab:Poisson}
\end{table}

\section{The symmetric-sector Dirichlet inverse}\label{gd:sec:spectral}

All operators in this section are evaluated at the finite shape $\psi_0$.
Set $e=a/2=\lambda$. Direct differentiation of $q^{-e}$ yields
\begin{equation}\label{gd:eq:conjugation}
 q^e\mathcal L(q^{-e}Z)=(\Delta_n+\mathcal V)Z,\qquad
 \mathcal V=\mu|p|^2-\frac{a(a-2)}2\frac{q_y}{yq}
                 -\frac{a(a-2)}4\frac{|\nabla q|^2}{q^2}.
\end{equation}
Here $\Delta_2=\partial_{x_1}^2+\partial_y^2$. The identity
$\Delta_2(yq)=0$ implies $\Delta_nq=(a-2)q_y/y$;
substitution cancels the first derivatives of $Z$.
Let $L=\Delta_n+\mathcal V$, $c=\rho b$, and $v_0=c^2$.

\subsection{The complete window}\label{gd:sec:window}
With reduced measure
$r^{n-1}\,dr\,(1-\xi^2)^{(n-3)/2}\,d\xi$, an orthonormal Dirichlet basis
in the invariant even sector is
\begin{equation}\label{gd:eq:spectral-basis}
 b_{l,\alpha}(r,\xi)=
 \frac{\sqrt2\,r^{-\lambda}J_{l+\lambda}(\alpha r)}
 {|J_{l+\lambda+1}(\alpha)|}
 \frac{C_l^\lambda(\xi)}{\sqrt{h_l^{(2)}}},\qquad
 h_l^{(2)}=\frac{\pi2^{1-2\lambda}\Gamma(l+2\lambda)}
 {l!(l+\lambda)\Gamma(\lambda)^2},
\end{equation}
where $l$ is even and $\alpha$ runs through all positive zeros of
$J_{l+\lambda}$. The Bessel equation and differentiation in $\alpha$
give the radial norm, as in~\eqref{r2:eq:basis}.
Angular orthogonality and the self-adjoint ball Dirichlet spectral
decomposition give completeness. The analytic weight $h_l$ in
\eqref{gd:eq:Y} is distinct from the $L^2$ normalisation $h_l^{(2)}$.

Let $\PP$ include every state with
\begin{equation}\label{gd:eq:window}
 |\alpha-c|<12,\qquad \QQ=I-\PP .
\end{equation}
There are two complete coverage arguments for this window.

\begin{lemma}\label{gd:lem:coverage}
Opposite-sign root brackets satisfying the spacing and guard conditions
below enumerate all states of~\eqref{gd:eq:window}.
\end{lemma}
\begin{proof}
Put $\nu=l+\lambda$. The leading series of $J_\nu$ makes
$J_\nu,J_\nu'>0$ near zero. At a first stationary point $x<\nu$
with $J_\nu>0$, the Bessel equation would give
$J_\nu''=(\nu^2/x^2-1)J_\nu>0$, contrary to a first loss of
positive derivative. A first zero of $J_\nu$ is also impossible
while its derivative is positive. Thus both remain positive for $x<\nu$.
If $J_\nu'(\nu)=0$, its equation gives $J_\nu''(\nu)=0$ and
$J_\nu'''(\nu)=-2J_\nu(\nu)/\nu<0$; Taylor expansion would make
$J_\nu'$ negative immediately to the left. Hence
$J_\nu,J_\nu'>0$ also at $\nu$.

For $z_\nu(x)=\sqrt{x}J_\nu(x)$ one has
\[
 z_\nu''+\left(1-\frac{\nu^2-1/4}{x^2}\right)z_\nu=0.
\]
On an interval ending at $B>\nu$, its coefficient is at most
$q_B=1-(\nu^2-1/4)/B^2$. The integration-by-parts and interval
Poincar\'e argument of Lemma~\ref{r2:lem:zero-coverage} gives spacing
at least $\pi/\sqrt{q_B}$. Starting from
$z_\nu(\nu),z_\nu'(\nu)>0$, comparison with the constant-coefficient
equation gives first-zero distance at least $\pi/(2\sqrt{q_B})$.
Consequently the first bracket ending strictly before
$\nu+3\pi/(2\sqrt{q_B})$ cannot miss an earlier zero.
If the outer endpoints of successive brackets have separation less
than $2\pi/\sqrt{q_B}$, no zero is missed between them.
Each bracket has opposite nonzero signs and width below the spacing;
it therefore contains exactly one simple zero.
Enumeration from $\nu$ through a guard root above $c+12$ covers the
window in that angular sector. Positivity up to $\nu$ excludes
all orders $\nu\ge c+12$.

For the second argument evaluate $J_\nu$ on the exact grid
$c-12+k/4$, $0\le k\le96$, in each remaining sector.
Here $\nu\ge3/2$, so the differential-equation coefficient is at
most one on the positive axis. Distinct positive roots are separated
by at least $\pi$. A grid cell contains at most one root.
Simplicity, and nonzero endpoint signs, make equal signs equivalent
to no root and opposite signs equivalent to one root.
Exclude equality at the two window edges. The same positivity bound excludes all larger angular orders.
\end{proof}

\subsection{The polynomial potential and its full error}\label{gd:sec:potential}
Index shape coefficients throughout this part by their actual odd degree:
$\psi_0(w)=\sum_{j\ {\rm odd}}c_jw^j$, $c_1=b$.
For a full, retained or discarded list define
\[
 s_1=\sum_{j>1}j|c_j|,\qquad
 s_2=\sum_{j>1}j(j-1)|c_j|,\qquad
 s_3=\frac13\sum_{j>1}j(j-1)(j-2)|c_j|.
\]
The integral formula for $q$ gives
$|p-b|,|q-b|\le s_1$, $|q_{x_1}|,|q_y|\le s_2$ and
$|q_y/y|\le s_3$. For the last bound, $q_y(x_1,0)=0$ because $q$ is even in $y$, and
$q_{yy}=-\int_0^1s^2\re p''(x_1+isy)\,ds$ gives
$|q_{yy}|\le s_3$, hence $|q_y|\le s_3|y|$.
Let the subscripts $f,t,d$ denote full, retained and discarded data.
Retain degrees at most $61$ in the polynomial potential only.
Put $q_*=b-s_{1,f}>0$ and $A_n=a(a-2)$. Then
\begin{equation}\label{gd:eq:vfull}
 \nrm{\mathcal V-v_0}_\infty\le
 v_f:=\mu(2bs_{1,f}+s_{1,f}^2)
 +\frac{A_ns_{3,f}}{2q_*}+\frac{A_ns_{2,f}^2}{2q_*^2}.
\end{equation}
The difference caused by discarding shape degrees is at most
\begin{equation}\label{gd:eq:shape-error}
 \begin{split}
 \delta_{\rm shape}={}&
 \mu(2b+s_{1,f}+s_{1,t})s_{1,d}\\
 &+\frac{A_n}2\left(\frac{s_{3,d}}{q_*}
                    +\frac{s_{3,t}s_{1,d}}{q_*^2}\right)\\
 &+\frac{A_n}4\left\{
 \frac{2(s_{2,f}+s_{2,t})s_{2,d}}{q_*^2}
 +\frac{2s_{2,t}^2(2b+s_{1,f}+s_{1,t})s_{1,d}}{q_*^4}\right\}.
 \end{split}
\end{equation}
The first line bounds the difference of $|p|^2$.
The second uses $|q_f^{-1}-q_t^{-1}|\le s_{1,d}/q_*^2$.
For the last line split the gradient square and its squared reciprocal;
$|\nabla q|^2\le2s_2^2$ and
$|q_f^{-2}-q_t^{-2}|\le
(2b+s_{1,f}+s_{1,t})s_{1,d}/q_*^4$ give the displayed bound.

Replace $q_t^{-1}$ by the polynomial
$S_4=b^{-1}\sum_{k=0}^4[-(q_t-b)/b]^k$ and $q_t^{-2}$ by $S_4^2$.
Since $s_{1,t}\le s_{1,f}<b$, the geometric remainder is at most
$\delta_S=(s_{1,t}/b)^5/(b-s_{1,t})$.
Also $|S_4|\le(b-s_{1,t})^{-1}$.
The resulting polynomial potential $\mathcal V_t$ satisfies
\begin{equation}\label{gd:eq:potential-error}
 \nrm{\mathcal V-\mathcal V_t}_\infty\le
 \delta_V:=\delta_{\rm shape}
 +\frac{A_n}2s_{3,t}\delta_S
 +\frac{A_ns_{2,t}^2\delta_S}{b-s_{1,t}}.
\end{equation}
This includes the squared-reciprocal error
$2\delta_S/(b-s_{1,t})$. As a polynomial in $(x_1,y)$,
$\mathcal V_t$ has degree at most $598$: the term
$|\nabla q_t|^2S_4^2$ has degree $118+480$.
Set $V_t=\mathcal V_t-v_0$ and $v_*=v_f+\delta_V$.

\subsection{The Gram identity and Schur complement}\label{gd:sec:Schur}
Form the exact finite matrices
\begin{equation}\label{gd:eq:matrices}
 \mathsf B=\PP V_t\PP,\qquad \mathsf B_2=\PP V_t^2\PP,\qquad
 H_{\mathsf P}=\diag(\alpha^2-v_0)-\mathsf B,\qquad G=\mathsf B_2-\mathsf B^2.
\end{equation}
The sans-serif symbols distinguish these matrices from the differential
operator $B$ in~\eqref{gd:eq:B}. The Gram identity is
\begin{equation}\label{gd:eq:Gram}
 G=(\QQ V_t\PP)^*(\QQ V_t\PP)\succeq0.
\end{equation}
It is the identity~\eqref{r2:eq:Gram} on the entire orthogonal complement;
in particular it includes every omitted radial and angular mode.

On the complement the ball operator $-\Delta_n-v_0$ has absolute
spectral gap at least $12(2c-12)$, with $c>12$.
Therefore
\[
 M_Q=\QQ(-\Delta_n-v_0-V_t)\QQ,\qquad
 \nrm{M_Q^{-1}}\le\gamma^{-1},\qquad
 \gamma=12(2c-12)-v_*>0 .
\]
This is the Neumann factorisation~\eqref{r2:eq:unbounded-factor};
the operator is allowed to be indefinite.

\begin{proposition}\label{gd:prop:Schur}
Let $C$ be a square real matrix of the window size, $D$ a real
diagonal matrix with nonzero diagonal, and suppose
\begin{equation}\label{gd:eq:Schur-gates}
 \begin{gathered}
 \nrm{C^TC-I}_2\le\epsilon_C<1,\\
 \nrm{|D|^{-1/2}(C^TH_{\mathsf P}C-D)|D|^{-1/2}}_2\le\epsilon_{\rm S},\\
 \gamma^{-1}\tr(|D|^{-1}C^TGC)\le\eta_{\rm S},\qquad\epsilon_{\rm S}+\eta_{\rm S}<1.
 \end{gathered}
\end{equation}
Set
\begin{equation}\label{gd:eq:L2}
 \begin{split}
 S_*&=\frac{1+\epsilon_C}{\min_i|D_i|(1-\epsilon_{\rm S}-\eta_{\rm S})},\\
 B_{\rm pol}&=\bigl(1+(v_*/\gamma)^2\bigr)S_*+\gamma^{-1},\\
 B_{L^2}&=\frac{B_{\rm pol}}{1-B_{\rm pol}\delta_V}\qquad(B_{\rm pol}\delta_V<1).
 \end{split}
\end{equation}
Then $L$ has a symmetric-sector Dirichlet inverse with
$\nrm{L_D^{-1}}_{L^2\to L^2}\le B_{L^2}$.
\end{proposition}
\begin{proof}
The proof of Proposition~\ref{r2:prop:Schur} applies with this ball basis
and $V_t$, the two defects being bounded directly in the $\ell^2$
operator norm. It uses only the orthogonal splitting of the Dirichlet
operator domain, the complement gap and~\eqref{gd:eq:Gram}.
In particular, for $T=\QQ V_t\PP C|D|^{-1/2}$,
\[
 \nrm{T^*M_Q^{-1}T}\le\gamma^{-1}\nrm{T^*T}
 \le\gamma^{-1}\tr(T^*T)\le\eta_{\rm S}.
\]
After congruence the Schur matrix is $\operatorname{sgn}(D)$ plus a
perturbation of norm at most $\epsilon_{\rm S}+\eta_{\rm S}$, so its inverse is bounded
by $S_*$. The block inverse has norm at most
$(1+(v_*/\gamma)^2)S_*+\gamma^{-1}$, exactly as in
\eqref{r2:eq:BLref}. Finally the bounded perturbation
$\mathcal V-\mathcal V_t$ is transferred by the factorisation of
\eqref{r2:eq:transfer}, giving $B_{L^2}$.
\end{proof}

The reduced measure in~\eqref{gd:eq:spectral-basis} differs from
Lebesgue measure on invariant functions on $B^n$ by the factor
$|S^{n-2}|$. On the even invariant sector, the map
$f(r,\xi)\mapsto |S^{n-2}|^{-1/2}f(|x|,x_1/|x|)$ is a unitary
identification of the reduced space with the invariant subspace of
$L^2(B^n,dx)$. It intertwines the Dirichlet operators, so $B_{L^2}$
is also the inverse bound on this subspace with Lebesgue measure.

\subsection{Quadrature enclosures}\label{gd:sec:quadrature}
Here are the finite expressions defining the matrix enclosures.
For odd $n$, the angular rule is Gauss--Legendre with polynomial weight
$(1-\xi^2)^{(n-3)/2}$.
For even $n$, it is Gauss--Chebyshev of the second kind, with
\[
 \xi_k=\cos\frac{k\pi}{N_\theta+1},\qquad
 \omega_k=\frac{\pi}{N_\theta+1}
             \sin^2\frac{k\pi}{N_\theta+1}\quad(1\le k\le N_\theta),
\]
and remaining polynomial weight $(1-\xi^2)^{(n-4)/2}$.
Choose the angular order so that the entire polynomial degree of each
integrand is at most $2N_\theta-1$; this makes angular integration exact.
This covers both integral and half-integral $\lambda$.

For radial Gauss--Legendre order $N_r$, the Bernstein ellipse of parameter
two in $2r-1$ gives $|r|\le9/8$ and $|\im r|\le3/8$.
Poisson's Bessel integral~\cite[10.9.4]{DLMF} gives
\[
 |z^{-\nu}J_\nu(z)|\le
 \frac{e^{|\im z|}}{2^\nu\Gamma(\nu+1)}.
\]
For each window basis vector its radial integrand factor, including
$r^{(n-1)/2}$, is bounded by
\begin{equation}\label{gd:eq:ellipse}
 \hat e_i=\frac{\sqrt2C_l^\lambda(1)}
 {\sqrt{h_l^{(2)}}|J_{\nu+1}(\alpha)|}
 \frac{\alpha^\nu(9/8)^{l+(n-1)/2}e^{3\alpha/8}}
 {2^\nu\Gamma(\nu+1)},\qquad \nu=l+\lambda .
\end{equation}
The products are entire: after removing the Bessel leading powers,
the radial power is $l+l'+n-1$, an integer.
Let $\hat e_*=\max_i\hat e_i$ and let $V_e$ bound $|V_t(r,\xi)|$ for $r$ on the ellipse
and $-1\le\xi\le1$, using the $s_i$ sums at $9/8$ and the finite
polynomial $S_4$. The entrywise radial errors can be taken as
\begin{equation}\label{gd:eq:quadrature-errors}
 q_1=32\hat e_*^2(1+V_e)2^{-2N_r},\qquad
 q_2=32\hat e_*^2(1+V_e^2)2^{-2N_r}.
\end{equation}
Indeed the Chebyshev coefficients of an analytic integrand bounded by
$M$ are at most $2M2^{-k}$. The radial Gaussian rule is exact through
degree $2N_r-1$; bounding both its positive mass and the integration
mass and summing the remaining geometric series gives the factor
$16$. The angular integration has mass at most two, giving $32$.

Let $u_{\rm fl}=2^{-52}$,
$\gamma_k=ku_{\rm fl}/(1-ku_{\rm fl})$, and $\tau=10^{-280}$.
Use these bounds only when $ku_{\rm fl}<1$, with finite intermediate
values. Regard every binary64 export as its exact dyadic rational.
At a tensor node the radial and angular factors are
\[
 \frac{\sqrt{2\omega_rr}J_\nu(\alpha r)}{|J_{\nu+1}(\alpha)|},
 \qquad
 C_l^\lambda(\xi)
 \sqrt{\omega_\theta w_{\rm poly}(\xi)/h_l^{(2)}}.
\]
Here $\omega_r$ is the radial Gauss--Legendre weight on $[0,1]$,
$\omega_\theta$ is the angular weight, and
$w_{\rm poly}(\xi)=(1-\xi^2)^{(n-3)/2}$ for odd $n$,
$(1-\xi^2)^{(n-4)/2}$ for even $n$.
Let their exported vectors be $r_i,a_i$, with enclosed norm errors
$e_{r,i},e_{a,i}$. The column error is bounded by
\begin{equation}\label{gd:eq:tensor-error}
 e_i^{\rm col}=e_{r,i}\nrm{a_i}_2+\nrm{r_i}_2e_{a,i}
              +e_{r,i}e_{a,i}
              +u_{\rm fl}\nrm{r_i}_2\nrm{a_i}_2+\tau .
\end{equation}
Thus $e_{\rm bas}=(\sum_i(e_i^{\rm col})^2)^{1/2}$
bounds the basis matrix Frobenius error.
Bessel argument uncertainty is bounded by differentiating the Poisson
integral:
\[
 \sup_{0\le x\le X}|J_\nu'(x)|
 \le\frac{\nu X^{\nu-1}+X^\nu}{2^\nu\Gamma(\nu+1)}
 \quad(\nu\ge1).
\]

The potential arrays are given by the exact identities
\begin{equation}\label{gd:eq:potential-arrays}
 \begin{aligned}
 q-b&=\sum_{j>1}c_jr^{j-1}U_{j-1}(\xi),&
 \re p-b&=\sum_{j>1}jc_jr^{j-1}T_{j-1}(\xi),\\
 q_{x_1}&=\sum_{j>1}jc_jr^{j-2}U_{j-2}(\xi),&
 q_y/y&=-2\sum_{j>1}c_jr^{j-3}C_{j-3}^{(2)}(\xi).
 \end{aligned}
\end{equation}
Each array is expanded in $r^iT_m(\xi)$ using
\[
 \begin{split}
 U_m&=\sum_{k=0}^{\lfloor m/2\rfloor}(2-\delta_{m,2k})T_{m-2k},\\
 C_N^{(2)}&=\sum_{k=0}^{\lfloor N/2\rfloor}
       (2-\delta_{N,2k})(k+1)(N-k+1)T_{N-2k}.
 \end{split}
\]
Here $S_{\rm pol}$ is the absolute sum of the resulting coefficients.
These coefficients are exported, and only the $T_m$ are evaluated by
the recurrence below.
For degree at most $d$, radial power error costs $\gamma_{d+1}$.
The recurrence at an exported node $|\xi|<1$ is
$\widehat T_{j+1}=2\xi\widehat T_j-\widehat T_{j-1}$,
$\widehat T_0=1$, $\widehat T_1=\xi$.
Suppose $10u_{\rm fl}(d+1)^3<10^{-3}$. If
$|\widehat T_i|\le1.001$ for $i\le j$, the product and subtraction
have total rounding error $|\epsilon_j|\le6u_{\rm fl}$, apart from
the separate underflow allowance. Since $|U_m|\le m+1$,
\[
 \left|\widehat T_{j+1}-T_{j+1}(\xi)\right|
 =\left|\sum_{i=1}^j\epsilon_iU_{j-i}(\xi)\right|
 \le3u_{\rm fl}j(j+1)<10^{-3}.
\]
Thus $|\widehat T_{j+1}|\le1.001$ by induction, and the error at
every degree $k\le d$ is at most
$3u_{\rm fl}d(d-1)\le10u_{\rm fl}(d+1)^3$, the recurrence allowance
of the array error.
Node-export uncertainty is separate: $|T_m|\le1$ and
$|T_m'|\le m^2$ imply the array error
$S_{\rm pol}(d\,\delta r+d^2\delta\xi)$.
Add coefficient export errors and both dense-product allowances.
Propagate them through the seven-variable polynomial potential by
interval differentiation on a box containing both the exact and
exported inputs, so the mean value theorem applies; an additional $2\gamma_{400}$ times the
absolute arithmetic majorant covers its fewer than $100$ scalar
operations. Let $\delta_{\rm pt}$ be the resulting pointwise potential error,
and $v_{\rm exp}$ the exported absolute bound.

For an exported basis matrix $E_F$ with $\nrm{E_F}_F\le K_F$,
window size $m$ and total node count $N_q$, valid matrix errors are
\begin{equation}\label{gd:eq:matrix-errors}
 \begin{split}
 e_B={}&v_*(2K_Fe_{\rm bas}+e_{\rm bas}^2)+\delta_{\rm pt} K_F^2
              +\gamma_{N_q+4}v_{\rm exp}K_F^2+mq_1+\tau,\\
 e_{B_2}={}&v_*^2(2K_Fe_{\rm bas}+e_{\rm bas}^2)
       +\delta_{\rm pt}(2v_*+\delta_{\rm pt})K_F^2
       +\gamma_{N_q+6}v_{\rm exp}^2K_F^2+mq_2+\tau .
 \end{split}
\end{equation}
These follow by adding and subtracting the exact factors in
$E_F^T\diag(V_F)E_F$. Include root-diagonal export, subtraction
and symmetrisation errors in $e_H$.
Bound $\epsilon_C$ by the Frobenius norm of the computed $C^TC-I$
plus $\gamma_{m+3}\nrm C_F^2+\tau$.
Let $H_F$ be the exported symmetrised $H_{\mathsf P}$. Bound the numerator of
$\epsilon_{\rm S}$ by the Frobenius norm of the computed $C^TH_FC-D$ plus
\[
 \gamma_{2m+8}(\nrm C_F^2\nrm{H_F}_F+\nrm D_F)+\tau
                    +(1+\epsilon_C)e_H.
\]
Division by $\min|D_i|$ bounds the normalised defect in
\eqref{gd:eq:Schur-gates}.
For the exported Gram difference $G_F=\mathsf B_{2,\rm exp}-\mathsf B_{\rm exp}^2$ add
\begin{equation}\label{gd:eq:Gram-error}
 e_G=e_{B_2}+e_B(2v_*+e_B)+\gamma_{m+4}\nrm{\mathsf B_{\rm exp}}_F^2
                  +u_{\rm fl}(\nrm{\mathsf B_{2,\rm exp}}_F+\nrm{\mathsf B_{\rm exp}^2}_F)+\tau .
\end{equation}
The weighted trace allowance is
\[
 \frac{\gamma_{4m+20}\nrm C_F^2\nrm{G_F}_F
              +m(1+\epsilon_C)e_G}{\min|D_i|}+\tau .
\]
Divide the upper trace bound by $\gamma$ to obtain $\eta_{\rm S}$.
For an exported array $A_F$ with $N$ entries, let $s_F$ be its computed
sum of squares. The Frobenius reduction uses
\[
 \nrm{A_F}_F\le\sqrt{\frac{s_F}{1-\gamma_{2N+2}}+\tau}.
\]

\section{Transfer of the Dirichlet inverse to the working norm}
\label{gd:sec:transfer}

\subsection{Low angular degrees by heat smoothing}\label{gd:sec:heat}
Let $M\ge\nrm{\mathcal V}_Y$, so $|\mathcal V|\le M$, and let $Lv=f$
with zero Dirichlet data in the invariant sector.
For the Dirichlet heat semigroup with this bounded real potential,
\begin{equation}\label{gd:eq:heat}
 \nrm{e^{tL}}_{2\to\infty}\le e^{Mt}(8\pi t)^{-n/4},
 \qquad
 \nrm{e^{tL}}_{\infty\to\infty}\le e^{Mt}\quad(t>0).
\end{equation}
To obtain these bounds, the maximum principle first dominates the
Dirichlet heat kernel by the positive free Gaussian kernel, after
extension of data by zero. Iterating Duhamel's formula for multiplication
by $\mathcal V$, taking absolute values, and using $|\mathcal V|\le M$
sums to the factor $e^{Mt}$.
The $L^2$ norm of
$(4\pi t)^{-n/2}\exp(-|x|^2/(4t))$ is $(8\pi t)^{-n/4}$.
Its $L^1$ norm is one. This proves both estimates, also on invariant
inputs with the full Lebesgue measure.

The spectral inverse places $v$ in the Dirichlet operator domain.
Integrating the semigroup derivative gives
\[
 v=e^{tL}v-\int_0^t e^{sL}f\,ds .
\]
Since $\nrm f_2\le\sqrt{|B^n|}\nrm f_Y$ and
$\nrm v_2\le B_{L^2}\nrm f_2$, it follows that
\begin{equation}\label{gd:eq:Bsup}
 \nrm v_\infty\le B_\infty\nrm f_Y,\qquad
 B_\infty=e^{Mt}(8\pi t)^{-n/4}\sqrt{|B^n|}B_{L^2}
                    +\frac{e^{Mt}-1}{M},
 \quad |B^n|=\frac{\pi^{n/2}}{\Gamma(n/2+1)}.
\end{equation}
We take $t=1/M$. In particular $v$ is bounded before any analytic
coefficient norm of $v$ is used.
The equation $\Delta_n v=f-\mathcal Vv$ now has bounded data;
full-ball elliptic regularity gives $v\in W^{2,p}$ for every finite $p$,
and hence a continuous radial expansion~\cite{GT01}.

The dimension of degree-$l$ spherical harmonics on $S^{n-1}$ is
\[
 d_l^{\rm sph}=\frac{2l+n-2}{n-2}\binom{l+n-3}{l}.
\]
The addition theorem gives
$\int_{S^{n-1}}G_l^2\,d\mathcal H^{n-1}/|S^{n-1}|=1/d_l^{\rm sph}$.
At each $r$, angular projection and Cauchy--Schwarz imply
\begin{equation}\label{gd:eq:low-angular}
 |v_l(r)|\le\sqrt{d_l^{\rm sph}}\nrm{v(r,\cdot)}_\infty
                 \le\sqrt{d_l^{\rm sph}}B_\infty\nrm f_Y .
\end{equation}
Only finitely many angular degrees will use this estimate.

\subsection{The entire high angular complement}\label{gd:sec:high}
Let $v_Y\ge\nrm{\mathcal V-v_0}_Y$.
The tabulated $L_0$ is the least even degree with $d_h>2v_Y$.
For the argument, choose any even $L_0$ with
\begin{equation}\label{gd:eq:high}
 d_h=L_0(L_0+n-2)-v_0>v_Y,\qquad \beta=v_Y/d_h<1 .
\end{equation}
For every $l\ge L_0$, the positive zero-order part of $-\mathcal D_l-v_0$
is at least $d_h$. The radial maximum principle therefore bounds its
zero-Dirichlet inverse in supremum norm by $1/d_h$.
First define this high-mode inverse in $L^2$ by coercivity.
For a source in $Y$, its finite angular truncations converge in $Y$;
the radial bound makes their inverse images converge in $Y$ and in
$L^2$. Closedness of the Dirichlet operator places their limit in
$H^2\cap H^1_0$ and identifies it with the same inverse.
Let $\Pi_H$ be angular projection to these degrees, and let
$v_{\rm low}=(I-\Pi_H)v$. By~\eqref{gd:eq:low-angular} this low part
belongs to $Y$. Solve
\[
 (\Delta_n+v_0)\tilde v_H
   =\Pi_H\bigl(f-(\mathcal V-v_0)(v_{\rm low}+\tilde v_H)\bigr)
\]
with zero Dirichlet data by contraction on the high part of $Y$.
The algebra bound gives contraction factor at most $\beta$.

We next identify $\tilde v_H$ with $\Pi_Hv$, as in the proof of
Proposition~\ref{r2:prop:Yinverse}.
Its Poisson right side, including $-v_0\tilde v_H$, lies in $Y$.
Thus $\tilde v_H$ equals its zero-Dirichlet Poisson potential, and belongs
to $H^2\cap H^1_0$. The original high projection belongs to the same
operator domain. Their difference $z$ has only degrees $l\ge L_0$ and
satisfies
\[
 (\Delta_n+v_0)z=-\Pi_H((\mathcal V-v_0)z).
\]
Angular orthogonality and $r\le1$ give
$\int|\nabla z|^2\ge L_0(L_0+n-2)\int|z|^2$.
Testing the equation with $-z$, using self-adjointness of $\Pi_H$,
gives
\[
 d_h\int|z|^2
 \le\int|\nabla z|^2-v_0\int|z|^2
 =\int(\mathcal V-v_0)|z|^2\le v_Y\int|z|^2 .
\]
Hence $z=0$. This proves membership of the original solution in $Y$
before rearranging any inequality involving its $Y$ norm.

Add the low norm to the high contraction inequality, use $h_l\le R^l$,
and obtain
\begin{equation}\label{gd:eq:BY}
 \nrm{L_D^{-1}f}_Y\le B_Y\nrm f_Y,\qquad
 B_Y=\frac{B_\infty
       \displaystyle\sum_{\substack{0\le l<L_0\\l\ {\rm even}}}
                R^l\sqrt{d_l^{\rm sph}}+d_h^{-1}}{1-\beta}.
\end{equation}
The Poisson trace formula applied to $\Delta_n v=f-\mathcal Vv$
also gives
\begin{equation}\label{gd:eq:trace-conjugated}
 \nrm{\partial_rv}_{W_1}\le(1+MB_Y)\nrm f_Y .
\end{equation}

\subsection{Undoing the conjugation}\label{gd:sec:unconjugation}
Let $\nrm{q-b}_Y\le q_b<b$ and define
\begin{equation}\label{gd:eq:source-multipliers}
 \mathfrak I=(b-q_b)^{-1},\qquad
 S_q=b^e(1-q_b/b)^{-|e-1|},\qquad
 M_0=b^{-e}(1-q_b/b)^{-e}.
\end{equation}
The binomial series gives
$\nrm{bq^{e-1}}_Y\le S_q$ and $\nrm{q^{-e}}_Y\le M_0$.
For a positive exponent, absolute falling-factorial coefficients are
bounded by the corresponding rising-factorial coefficients.
This proves the stated enlarged bound for either integral or
half-integral $e$.

For the original Dirichlet equation $B\Phi=f$, set $Z=q^e\Phi$.
Then $LZ=bq^{e-1}f$. Write
\[
 \begin{gathered}
 X_*\ge\nrm{q_{x_1}}_{\mathcal Y_n},\quad Y_*\ge\nrm{q_y/y}_Y,\quad
 Y_2\ge\nrm{y^2}_Y,\quad Q^{[1]}\ge\mathfrak m_1(q-b),\\
 B_Z=S_qB_Y,\qquad J_Z=S_q+MB_Z .
 \end{gathered}
\]
Consequently $\nrm Z_Y\le B_Z\nrm f_Y$ and
$\nrm{\Delta_nZ}_Y\le J_Z\nrm f_Y$.
Define
\begin{equation}\label{gd:eq:JW}
 \begin{split}
 J_W=M_0\bigl[&
 J_Z+2e\mathfrak I(X_*C_x^K+Y_*C_y^K)J_Z\\
 &+\{e(n-4)\mathfrak I Y_*
       +e(e+1)\mathfrak I^2(X_*^2+Y_2Y_*^2)\}B_Z\bigr],\\
 N_W&=M_0(1+e\mathfrak I Q^{[1]})J_Z .
 \end{split}
\end{equation}

\begin{proposition}\label{gd:prop:original-inverse}
The original symmetric-sector zero-Dirichlet inverse satisfies
\[
 \nrm{\Delta_nB_D^{-1}f}_Y\le J_W\nrm f_Y,\qquad
 \nrm{\partial_rB_D^{-1}f}_{W_1}\le N_W\nrm f_Y .
\]
\end{proposition}
\begin{proof}
Expand
\[
 \begin{split}
 \Delta_n(q^{-e}Z)=q^{-e}\bigl[
 \Delta_nZ-2e q^{-1}\nabla q\cdot\nabla Z
 +\{-e q^{-1}\Delta_nq
       +e(e+1)q^{-2}|\nabla q|^2\}Z\bigr].
 \end{split}
\]
The zero-Dirichlet identity $Z=K_n^D\Delta_nZ$ and
Lemma~\ref{gd:lem:Poisson} control both gradient terms.
Use $\Delta_nq=(n-4)q_y/y$ and
$|\nabla q|^2=q_{x_1}^2+y^2(q_y/y)^2$ to obtain $J_W$.
At the boundary $Z=0$, so $\partial_r\Phi=q^{-e}\partial_rZ$.
Lemma~\ref{gd:lem:Poisson} bounds $\partial_rZ$ by $J_Z\nrm f_Y$
in $W_1$. Moreover
\[
 \nrm{q^{-e}|_\partial}_{W_1}
 \le M_0+eM_0\mathfrak I\nrm{\partial_\theta q|_\partial}_{W_0}
 \le M_0(1+e\mathfrak I Q^{[1]}).
\]
The boundary algebra of~\eqref{r2:eq:Wone} gives $N_W$.
\end{proof}

\begin{table}[htbp]
\centering\small
\def\gdRow#1#2#3#4#5#6#7#8{$#1$ & #4\\}
\begin{tabular}{@{}crrrrr@{}}
\toprule
$n$ & states & $B_{L^2}$ & $B_Y$ & $J_W$ & $N_W$\\
\midrule
\gdData
\bottomrule
\end{tabular}
\caption{Complete Dirichlet inverse bounds. $B_Y$ is for the conjugated
operator; $J_W,N_W$ are for the original operator $B$.}
\label{gd:tab:inverse}
\end{table}

\section{The exactly clamped centre and its residual}\label{gd:sec:centre}

\subsection{Finite field and infinite clamping lift}
\label{gd:sec:clamp}
The finite input arrays are exact decimals $c_{2j+1},a_j,s_j$,
$0\le j\le J$, with $s_j>0$: these are the keys \texttt{c},
\texttt{a}, \texttt{scales} of the frozen centre, and $J+1$ is the
mode count of Table~\ref{gd:tab:inputs}. Their shape convention is
$\psi_0(w)=\sum_{j\ {\rm odd}}c_jw^j$, $b=c_1$ and $p_0=\psi_0'$.
Put, in unscaled meridian coordinates $(X_1,Y_\psi)$,
\begin{equation}\label{gd:eq:physical-field}
 u_*=\sum_{j=0}^J\frac{a_j}{s_j}
 \mathcal R^{-\lambda}J_{2j+\lambda}(\rho\mathcal R)
 G_{2j}(X_1/\mathcal R),\qquad
 \mathcal R^2=X_1^2+Y_\psi^2,\qquad U_*=u_*\circ\Theta_{\psi_0}.
\end{equation}
The Bessel power series expresses each summand as a solid harmonic
times an entire series in $\mathcal R^2$, and separation gives
$(\Delta_n+\mu)u_*=0$. The $s_j$ are fixed divisors; no identity
between them and Bessel values is required. By~\eqref{gd:eq:physical-multiplier},
$BU_*=0$ exactly.

Let $D_l,N_l$ be the full Gegenbauer coefficients of
$D=U_*|_{r=1}-1$ and $N=\partial_rU_*|_{r=1}$.
The double-trace lifts are
\begin{equation}\label{gd:eq:lifts}
 h_l^D=r^l[1-l(r^2-1)/2]G_l,\qquad
 h_l^N=\tfrac12r^l(r^2-1)G_l .
\end{equation}
Their value and radial derivative traces are respectively $(G_l,0)$
and $(0,G_l)$.
Define the exact centre by
\begin{equation}\label{gd:eq:clamp}
 \ell=\sum_l(D_lh_l^D+N_lh_l^N),\qquad
 U_0=U_*-\ell,\qquad g_0=\Delta_nU_0,\qquad x_0=(g_0,p_0).
\end{equation}
The boundary decay proved below gives analytic convergence on a
neighbourhood of the closed ball. Thus $U_0=1$ and $\partial_rU_0=0$
on the sphere, $g_0\in\Gc$, and $U_0=1+K_ng_0$ by Dirichlet uniqueness.
This is the clamping construction of~\eqref{r2:eq:clamped}
with solid zonal harmonics in place of planar harmonics.

\subsection{Complex majorants for the finite field}\label{gd:sec:strip}
Choose $S>R>1$, $\delta>0$, and $r_0=1/\rho<1$.
For $0\le r\le1$, $|\im\theta|\le\log S$, allow independent
perturbations of size $\delta$ in $w=re^{i\theta}$ and
$\bar w=re^{-i\theta}$. Define
\begin{equation}\label{gd:eq:strip-parameters}
 \begin{split}
 s&=b\delta S/r_0+
       S\sum_{j>1}|c_j|(S+\delta/r_0)^j,\qquad d_s=2s+s^2/b,\\
 z_*&=\sum_{j\ {\rm odd}}|c_j|(r_0S+\delta)^j,\qquad
 H_s=S(b+s),\qquad A_s=H_s/(b-d_s).
 \end{split}
\end{equation}
Assume $d_s<b$.
For $r\ge r_0$ write the two perturbed coordinates as
$X_1\pm iY_\psi=re^{\pm i\theta}A_\pm$.
Then $|A_\pm-b|\le s$.
To verify the nonlinear part of this estimate, set
$\varsigma=|re^{\pm i\theta}|\in[r_0/S,S]$.
For odd $m\ge3$, $f(\varsigma)=(\varsigma+\delta)^m/\varsigma$ has derivative
$f'(\varsigma)=(\varsigma+\delta)^{m-1}((m-1)\varsigma-\delta)/\varsigma^2$, so $f$ decreases for
$\varsigma<\delta/(m-1)$ and increases afterwards; hence its supremum
on $[r_0/S,S]$ is attained at an endpoint. Since $r_0<1<S$,
\[
 \begin{split}
 f(S)&=(S+\delta)^m/S\le S(S+\delta/r_0)^m,\\
 f(r_0/S)&=Sr_0^{m-1}(1/S+\delta/r_0)^m
                       \le S(S+\delta/r_0)^m.
 \end{split}
\]
The linear perturbation is at most $b\delta S/r_0$.

The product $A_+A_-$ lies in the disc of radius $2bs+s^2=bd_s<b^2$
about $b^2$, so it admits a square root $B_s$ with positive real part.
It satisfies $|B_s-b|\le d_s$ because $|B_s+b|\ge b$.
The complex radius is $rB_s$, with positive real part
and $|\im(rB_s)|\le d_s$.
The positive cosine expansion of the solid harmonic gives
$|H_l(X_1,Y_\psi)|\le(rH_s)^l$ and
$|G_l(X_1/\mathcal R)|\le A_s^l$.

For $\nu>0$, $\re z>0$, Schl\"afli's integral~\cite[10.9.6]{DLMF} gives
\[
 |J_\nu(z)|\le e^{|\im z|}+\frac1{\pi\nu}.
\]
Indeed the finite cosine integral has modulus at most $e^{|\im z|}$,
and the remaining integral is bounded by
$\pi^{-1}\int_0^\infty e^{-\nu t}\,dt$.
The normalised Poisson integral also applies at the removable origin.
Writing $l=2j$, $\nu=l+\lambda$ and
$P_\nu=(\rho/2)^\nu/\Gamma(\nu+1)$, a uniform field majorant is
\begin{equation}\label{gd:eq:Mfield}
 \begin{split}
 M_*=\sum_j\frac{|a_j|}{s_j}\biggl[
 P_\nu z_*^l e^{\rho z_*}
 +\min\biggl\{&
 (r_0(b-d_s))^{-\lambda}A_s^l
                 \left(e^{\rho d_s}+\frac1{\pi\nu}\right),\\[-2pt]
 &P_\nu H_s^l e^{\rho d_s}\biggr\}\biggr].
 \end{split}
\end{equation}
The first term covers $r\le r_0$ by the entire solid-harmonic
representation; the minimum covers $r\ge r_0$ by the two Bessel bounds.
Adding these nonnegative bounds enlarges their maximum.

Angular Cauchy estimates and Lemma~\ref{gd:lem:connection} give, with
$C=\kappa(S+R)/(S-R)$,
\begin{equation}\label{gd:eq:fixed-field-bounds}
 \begin{gathered}
 U_b=CM_*,\quad X_b=2U_b/\delta,\quad Y_b=RX_b,\\
 XX_b=6U_b/\delta^2,\quad YX_b=RXX_b,\quad YY_b=R^2XX_b .
 \end{gathered}
\end{equation}
These bound $U_*,(U_*)_{x_1},y(U_*)_y,(U_*)_{x_1x_1},
y(U_*)_{x_1y},y^2(U_*)_{yy}$ in $Y$ or $\mathcal Y_n$
according to parity.
For the first derivative the two independent complex coordinates give
the factor two; their second derivatives and twice the mixed derivative
give six by the polydisc Cauchy estimates.
The factors $y,y^2$ are multiplied in the raw Fourier algebra before
conversion of the resulting regular scalar. Thus $y$ itself is not
treated as an analytic scalar on $B^n$.
For these fixed-field estimates use $S=(102/100)R$ and
$\delta=(1000\rho)^{-1}$.
The boundary moment estimates use their own larger strip radius $S$.

\subsection{Boundary coefficients and their entire tails}\label{gd:sec:boundary}
Let $K$ and $J_b<K$ be the boundary sample and retained coefficient
parameters. Reflection of the values at
$\theta_k=(k+1/2)\pi/(2K)$, $0\le k<K$, gives $2K$
equispaced samples in $\phi=2\theta$.
For $D$ and $N$ the complex strip majorants are
$M_D=M_*+1$, $M_N=2SM_*/\delta$.
The corresponding continuum cosine coefficient of frequency $2j$
has modulus at most $2M_{D,N}S^{-2j}$.
The exact discrete Fourier transform, with its half-grid phase,
differs at each retained index $0\le j\le J_b$ by at most
\begin{equation}\label{gd:eq:alias}
 e_Q=\frac{4M_{D,N}S^{-2(2K-J_b)}}{1-S^{-4K}} .
\end{equation}
This is the alias sum at indices $j+2Kq$, as in
\eqref{r2:eq:alias}; the constant mode is included with the enlarged
same bound. Above $J_b$ use the full continuum coefficient bound.

Define the actual Gegenbauer moments
\begin{equation}\label{gd:eq:boundary-moments}
 D^{[k]}=\sum_lh_l(l+1)^k|D_l|,\qquad
 N^{[k]}=\sum_lh_l(l+1)^k|N_l|,\qquad 0\le k\le3.
\end{equation}
The connection lowers degree and hence the moment factor. Below,
$D^{[k]},N^{[k]}$ denote the upper bounds obtained by multiplying
the retained cosine moments
$\sum_{j=0}^{J_b}R^{2j}(2j+1)^k(|d_j|+e_Q)$ and the full tail
of~\eqref{gd:eq:tail-sums} by $\kappa$, where $d_j$ denotes the
corresponding cosine coefficient of $D$ or $N$. All later constants
are monotone in these bounds. For $q_{\rm tail}=(R/S)^2$, expand
$(2(J_b+1+t)+1)^k$ and use
\begin{equation}\label{gd:eq:tail-sums}
 \begin{aligned}
 \sum_{t\ge0}q_{\rm tail}^t&=\frac1{1-q_{\rm tail}},&
 \sum_{t\ge0}tq_{\rm tail}^t&=\frac{q_{\rm tail}}{(1-q_{\rm tail})^2},\\
 \sum_{t\ge0}t^2q_{\rm tail}^t&=\frac{q_{\rm tail}(1+q_{\rm tail})}{(1-q_{\rm tail})^3},&
 \sum_{t\ge0}t^3q_{\rm tail}^t&=\frac{q_{\rm tail}(1+4q_{\rm tail}+q_{\rm tail}^2)}{(1-q_{\rm tail})^4}.
 \end{aligned}
\end{equation}
The geometric coefficient decay also shows
that the lift in~\eqref{gd:eq:clamp} converges with all its derivatives
on a neighbourhood of the closed ball.

Direct differentiation of~\eqref{gd:eq:lifts} gives
\[
 \Delta_n h_l^D=-l(2l+n)H_l,\qquad
 \Delta_n h_l^N=(2l+n)H_l .
\]
Consequently set
\begin{equation}\label{gd:eq:lift-moments}
 \Lambda_k=2^k(D^{[k+1]}+N^{[k]})\ (0\le k\le2),\qquad
 \Lambda^\Delta_k=(n+2)(D^{[k+2]}+N^{[k+1]})\ (0\le k\le1).
\end{equation}
These bound the degree moments $\mathfrak m_k(\ell)$ and
$\mathfrak m_k(\Delta_n\ell)$.
For the first use that a lift monomial has degree at most
$l+2\le2(l+1)$, with total absolute coefficient at most $l+1$
for $h_l^D$ and one for $h_l^N$.
For the second use $2l+n\le(n+2)(l+1)$.
The radial factors themselves also give
\[
 \nrm\ell_Y\le D^{[0]}+N^{[0]}/2:
\]
the Dirichlet lift increases to one on $[0,1]$, and the Neumann
lift has modulus at most $1/2$.

\subsection{Residual and centre derivatives}\label{gd:sec:residual}
Let $Q_*,X_*,Y_*,P_*$ bound $q,q_{x_1},q_y/y,|p_0|^2$
in $\mathcal Y_n$, and let $Q^{[1]},X^{[1]},Y^{[1]},P^{[1]}$ bound their
total-degree coefficient moments.
Here $X_*,Y_*$ are the bounds of Section~\ref{gd:sec:unconjugation},
and $Q_*=b+q_b$. They arise from the finite shape identities~\eqref{gd:eq:potential-arrays}
and the full product
\[
 |p_0|^2=\sum_{i,j}(2i+1)(2j+1)c_{2i+1}c_{2j+1}
                         r^{2(i+j)}\cos(2(i-j)\theta).
\]
Retained connection coefficients can be converted directly;
Lemma~\ref{gd:lem:connection} bounds every higher angular degree.
Products add degree moments with nonnegative majorants.

Since $BU_*=0$, the actual residual is $E=F(x_0)=-B\ell$.
The product and derivative bounds give
\begin{equation}\label{gd:eq:E0}
 E_0=\frac{Q_*\Lambda^\Delta_0+
       a(X_*C_x^{\rm pol}+Y_*C_y^{\rm pol})\Lambda_1+\mu Q_*P_*\Lambda_0}{b}
       \ge\nrm E_Y
\end{equation}
and the coefficient degree moment bound
\begin{equation}\label{gd:eq:Ed}
 \begin{split}
 E_d=\frac1b\bigl[&
 Q^{[1]}\Lambda^\Delta_0+Q_*\Lambda^\Delta_1\\
 &+a\{(X^{[1]}C_x^{\rm pol}+Y^{[1]}C_y^{\rm pol})\Lambda_1
       +(X_*C_x^{\rm pol}+Y_*C_y^{\rm pol})\Lambda_2\}\\
 &+\mu\{(Q^{[1]}P_*+Q_*P^{[1]})\Lambda_0+Q_*P_*\Lambda_1\}\bigr].
 \end{split}
\end{equation}
Thus
\[
 \nrm{E_{x_1}}_{\mathcal Y_n}\le C_x^{\rm pol}E_d,\quad
 \nrm{yE_y}_Y\le C_y^{\rm pol}E_d,\quad
 \nrm{\partial_\theta E|_\partial}_{W_0}\le E_d .
\]

The inverse below uses derivatives of the fixed centre only.
Define $A^q=q_{x_1}/q$, $B^q=(q_y/y)/q$, $C^q=\mu|p_0|^2$.
With $\mathfrak I$ of~\eqref{gd:eq:source-multipliers}, upper bounds for these
coefficients and their derivatives are
\begin{equation}\label{gd:eq:coefficient-derivatives}
 \begin{aligned}
 A_*^q&=X_*\mathfrak I,& B_*^q&=Y_*\mathfrak I,& C_*^q&=\mu P_*,\\
 A_x^q&=C_x^{\rm pol}X^{[1]}\mathfrak I+X_*^2\mathfrak I^2,&
 A_y^q&=C_y^{\rm pol}X^{[1]}\mathfrak I+X_*Y_2Y_*\mathfrak I^2,\\
 B_x^q&=C_x^{\rm pol}Y^{[1]}\mathfrak I+Y_*X_*\mathfrak I^2,&
 B_y^q&=C_y^{\rm pol}Y^{[1]}\mathfrak I+Y_*^2Y_2\mathfrak I^2,\\
 C_x^q&=\mu C_x^{\rm pol}P^{[1]},&
 C_y^q&=\mu C_y^{\rm pol}P^{[1]} .
 \end{aligned}
\end{equation}
Here subscripts $x$ and $y$ indicate $\partial_{x_1}$ and
$y\partial_y$ respectively.
Use $\Delta_nU_*=-a(A^q(U_*)_{x_1}+B^qy(U_*)_y)-C^qU_*$
and~\eqref{gd:eq:fixed-field-bounds} to define
\begin{equation}\label{gd:eq:gstar}
 \begin{split}
 \Gamma_b={}&a(A_*^qX_b+B_*^qY_b)+C_*^qU_b,\\
 \Gamma_x={}&a(A_x^qX_b+A_*^qXX_b+B_x^qY_b+B_*^qYX_b)
                      +C_x^qU_b+C_*^qX_b,\\
 \Gamma_y={}&a(A_y^qX_b+A_*^qYX_b+B_y^qY_b+B_*^q(Y_b+YY_b))
                      +C_y^qU_b+C_*^qY_b .
 \end{split}
\end{equation}
The term $Y_b+YY_b$ is the derivative of $y(U_*)_y$ by
$y\partial_y$. Finally the bounds for the clamped centre are
\begin{equation}\label{gd:eq:centre-bounds}
 \begin{gathered}
 U_c=U_b+D^{[0]}+N^{[0]}/2,\quad
 U_x=X_b+C_x^{\rm pol}\Lambda_1,\quad U_y=Y_b+C_y^{\rm pol}\Lambda_1,\\
 \Gamma_0=\Gamma_b+\Lambda^\Delta_0,\quad \Gamma_{0x}=\Gamma_x+C_x^{\rm pol}\Lambda^\Delta_1,\quad
 \Gamma_{0y}=\Gamma_y+C_y^{\rm pol}\Lambda^\Delta_1 .
 \end{gathered}
\end{equation}
They bound $U_0,(U_0)_{x_1},y(U_0)_y,g_0,
(g_0)_{x_1},y(g_0)_y$ in their respective parity spaces.

\section{The exact inverse of the corrected linearisation}
\label{gd:sec:coupled}

\subsection{The material identity, including the axis term}
\label{gd:sec:material}
All coefficients here are evaluated at $x_0$.
For $\delta p\in W$ define the odd real-coefficient holomorphic functions
\[
 \eta(w)=\int_0^w\delta p(z)\,dz,\qquad v=\eta/p_0 .
\]
The associated full-ball vector field is
$(\re v,(\im v/y)x')$, regular on the axis.
Its action on $U_0$ is $T=\re v\,(U_0)_{x_1}+\im v\,(U_0)_y$.

To derive the material identity, let $\Theta_t$ vary the shape,
$m_t=q_t|p_t|^2/b$, and keep the function
$u_0=U_0\circ\Theta_0^{-1}$ fixed.
The centre is analytic on a neighbourhood of the closed ball, so this
function is defined on a neighbourhood as well.
Equivalently the following differentiation is a local differential
identity. Differentiating
$B_t(u_0\circ\Theta_t)=m_t((\Delta+\mu)u_0)\circ\Theta_t$ gives
\begin{equation}\label{gd:eq:material-general}
 BT+(\delta B)U_0
 =v\cdot\nabla E+(\delta\log m-v\cdot\nabla\log m)E .
\end{equation}
For example the last transport term before expansion is
$m\,v\cdot\nabla(E/m)$, which explains the minus sign.

The cancellation for $|p|^2$ uses $\eta'=p'v+pv'$:
\[
 \delta\log|p|^2-v\cdot\nabla\log|p|^2=2\re v'.
\]
For $q$ one has, with $Y_\psi=yq=\im\psi$,
$v\cdot\nabla Y_\psi=\im(pv)=\im\eta$.
Hence $v\cdot\nabla q=\im\eta/y-q\im v/y$ and
\[
 \delta\log q-v\cdot\nabla\log q=\im v/y .
\]
Define
\begin{equation}\label{gd:eq:correction}
 \Ecal(\delta p)=\re v\,E_{x_1}+\im v\,E_y
                         +(2\re v'+\im v/y)E,\qquad
 \widetilde D=DF(x_0)-\Ecal .
\end{equation}
The coefficient of $\im v/y$ is one, independent of $n$, because
the multiplier in~\eqref{gd:eq:physical-multiplier} contains one factor
$q$. The dimension enters the operator $B$ through $a=n-2$.
For the divergence multiplier $|p|^2q^a$, the same calculation gives
coefficient $a=n-2$ instead.
Equations~\eqref{gd:eq:F} and~\eqref{gd:eq:material-general} now imply
\begin{equation}\label{gd:eq:corrected-factor}
 \widetilde D(\delta g,\delta p)=B(K_n\delta g-T).
\end{equation}
The material mechanism is the Hadamard shape derivative used in
Proposition~\ref{r2:prop:coupled}; see, for example,~\cite{SZ92,HP18}.

\subsection{Boundary reconstruction and both inverse identities}
\label{gd:sec:inverse-identities}
Exact clamping gives $\nabla U_0=0$ on the sphere.
Differentiating this identity in each tangential direction shows that
the Hessian annihilates every tangent vector. By symmetry of the
Hessian,
\[
 D^2U_0=(\Delta_nU_0)\mathbf n\otimes\mathbf n\quad\hbox{on }\partial B^n .
\]
Thus, with $w=e^{i\theta}$ on the meridian circle,
\begin{equation}\label{gd:eq:boundary-material}
  \begin{gathered}
 T=0,\qquad \partial_rT=-d_\partial\re(v\bar w),\\
 d_\partial=-\Delta_nU_0|_\partial=\mu|p_0|^2-bE/q .
 \end{gathered}
\end{equation}
The term from differentiating $v$ multiplies $\nabla U_0$ and vanishes.
At the poles these identities are full-ball identities, hence also
hold by continuity.

Assume that $B_D^{-1}$ exists as proved above, $p_0^{-1}$ belongs
to the holomorphic weighted algebra, and $d_\partial^{-1}\in W_1$.
For $f\in Y$ solve
\begin{equation}\label{gd:eq:Dirichlet-recipe}
 B\Phi=f,\quad \Phi|_\partial=0,\qquad N_\partial=\partial_r\Phi,\quad H_\partial=N_\partial/d_\partial.
\end{equation}
Write its cosine coefficients as
$H_\partial=H_{\partial,0}+\sum_{j>0}H_{\partial,2j}\cos(2j\theta)$ and set
\begin{equation}\label{gd:eq:realpart-recipe}
 \mathfrak a(w)=H_{\partial,0}+\sum_{j>0}H_{\partial,2j}w^{2j},\qquad v=w\mathfrak a.
\end{equation}
Then $\re(v\bar w)=H_\partial$ on the circle and
$\nrm{v'}_W=\nrm{\mathfrak a}_{W_1}=\nrm{H_\partial}_{W_1}$.
The holomorphic real-part problem is the one of
\eqref{r2:eq:realpart}, with cosine coefficients in place of paired
complex Fourier coefficients. Its real-coefficient normalisation
excludes the imaginary constant; the real constant $H_{\partial,0}$ determines
the dilation variation.

\begin{proposition}[Two-sided inverse]\label{gd:prop:coupled}
Under these hypotheses the corrected derivative
$\widetilde D:X\to Y$ has a bounded two-sided inverse $A:Y\to X$.
It is given by
\begin{equation}\label{gd:eq:coupled-recipe}
 \delta\psi=p_0v,\quad \delta p=(p_0v)',\quad
 \delta U=\Phi+T,\quad \delta g=\Delta_n\delta U,\qquad
 Af=(\delta g,\delta p).
\end{equation}
\end{proposition}
\begin{proof}
Equation~\eqref{gd:eq:boundary-material} makes both traces of
$\delta U$ zero, since $\partial_r\delta U=N_\partial-d_\partial H_\partial=0$.
The Laplacian estimate in the next subsection proves
$\delta g\in Y$; the moment formula therefore gives
$\delta g\in\Gc$ and Dirichlet uniqueness gives
$K_n\delta g=\delta U$. The shape variation is even holomorphic
with real coefficients.
Then~\eqref{gd:eq:corrected-factor} gives $\widetilde D Af=f$.

For the other identity, take any $(\delta g,\delta p)\in X$,
form its $\eta,v,T$, and put $\Phi=K_n\delta g-T$.
It has zero Dirichlet trace and satisfies
$B\Phi=\widetilde D(\delta g,\delta p)$.
Dirichlet uniqueness recovers precisely this $\Phi$ when the recipe is
applied to that right side.
Its normal trace is $N_\partial=d_\partial\,\re(v\bar w)$.
The function $\mathfrak a=v/w=\eta/(wp_0)$ has a removable singularity at zero.
It belongs to $W_1$, since
$\nrm{\eta/w}_{W_1}=\nrm{\delta p}_W$ and
$p_0^{-1}\in W_1$ by differentiation of its reciprocal.
The unique real-coefficient real-part solution therefore recovers
$\mathfrak a$ and then $v,\delta p,\delta U,\delta g$.
Hence $A\widetilde D=I_X$ as well as $\widetilde DA=I_Y$.
The bounds below prove boundedness on the entire spaces.
\end{proof}

The sector is essential to the Dirichlet inverse.
For a Schiffer solution, physical translation derivatives solve the
homogeneous Dirichlet equation. The axial derivative is odd in the
first coordinate and the transverse derivatives are not invariant
under $O(n-1)$, so they lie outside this sector.

\subsection{Cancellation of higher shape derivatives}
\label{gd:sec:no-loss}
Cauchy--Riemann cancellation gives
\begin{equation}\label{gd:eq:DeltaT}
 \begin{split}
 \Delta_nT={}&v\cdot\nabla g_0+2\re v'\,g_0
       -a\frac{\im v'}y(U_0)_{x_1}\\
       &+a\frac{\im v/y-\re v'}{y^2}\,y(U_0)_y .
 \end{split}
\end{equation}
To see the axis terms directly, write $A_v=\re v$, $C_v=\im v$.
The planar Laplacian gives
$\Delta_2(v\cdot\nabla U_0)
=v\cdot\nabla\Delta_2U_0+2\re v'\Delta_2U_0$.
For the remaining operator $ay^{-1}\partial_y$, subtracting its
transport derivative and $2\re v'$ times that operator leaves
\[
 a\frac{(A_v)_y}{y}(U_0)_{x_1}
 +a\left(\frac{(C_v)_y-2\re v'}y+\frac{C_v}{y^2}\right)(U_0)_y .
\]
Use $(A_v)_y=-\im v'$ and $(C_v)_y=\re v'$ to obtain
\eqref{gd:eq:DeltaT}. Thus no second derivative of the unknown
holomorphic variation appears.

For the odd monomial $v=w^j$, $j\ge3$, the regular quotient is
\[
 \frac{\im v/y-\re v'}{y^2}
                  =2r^{j-3}C_{j-3}^{(2)}(\xi);
\]
for $j=1$, both $\im v'/y$ and
$(\im v/y-\re v')/y^2$ vanish. The positive cosine expansions of $U_m$
and $C_m^{(2)}$, summed with weight $R$, give
\begin{equation}\label{gd:eq:axis-multipliers}
 \begin{gathered}
 c_a=\frac2{R(1-R^{-2})},\qquad
 c_v=\frac4{R^2(1-R^{-2})^2},\\
 \nrm{\im v'/y}_{\rm raw}\le c_a\nrm{v'}_W,\qquad
 \nrm{(\im v/y-\re v')/y^2}_{\rm raw}\le c_v\nrm{v'}_W,\\
 \nrm{\im v/y}_{\rm raw}\le\nrm{v'}_W,\qquad
 \nrm{\re v}_{\rm raw}\le R\nrm{v'}_W .
 \end{gathered}
\end{equation}
For the first inequality use the sum of the coefficients of
$U_m$ with decreasing powers $R^{-2k}$, bounded by
$2/(1-R^{-2})$ after the leading power.
For the second use the Gegenbauer generating coefficients
$(2)_k/k!=k+1$; their weighted sum is $(1-R^{-2})^{-2}$.
The leading powers and division by the derivative coefficient $j$
give the stated enlarged $c_a,c_v$.
For $\im v/y$ the unweighted sum of the coefficients of $U_{j-1}$
is $j$, and every frequency is at most $j-1$.
The last inequality follows coefficientwise from
$R^j/j\le R R^{j-1}$.
Lemma~\ref{gd:lem:connection} converts each regular scalar once.
Consequently
\begin{equation}\label{gd:eq:MT}
 M_T=\kappa\left[R\Gamma_{0x}+\Gamma_{0y}+2\Gamma_0+a c_aU_x+a c_vU_y\right]
\end{equation}
bounds $\nrm{\Delta_nT}_Y$ per unit $\nrm{v'}_W$.
This proves the field membership required in
Proposition~\ref{gd:prop:coupled} without changing the analytic radius.

\subsection{Reciprocals, inverse norm and left Newton defect}
\label{gd:sec:inverse-bounds}
Define the finite shape quantities
\[
 P=\nrm{p_0}_W,\qquad
 P'=\sum_{m>0}m|(p_0)_m|R^{m-1},\qquad q_0=2b-P>0 .
\]
The holomorphic reciprocal of $p_0$ and the raw scalar reciprocal
of $q$ have norm at most $1/q_0$, by Neumann series about $b$.
Set
\begin{equation}\label{gd:eq:boundary-reciprocal}
 \begin{gathered}
 d_*=\mu(2b^2-P^2)-bE_0/q_0>0,\\
 D_\theta=2\mu PRP'+b(E_d/q_0+E_0RP'/q_0^2),\\
 I_0=d_*^{-1},\qquad I_1=I_0+I_0^2D_\theta .
 \end{gathered}
\end{equation}
Indeed
$\nrm{d_\partial-\mu b^2}_{W_0}\le\mu(P^2-b^2)+bE_0/q_0<\mu b^2$,
so $\nrm{d_\partial^{-1}}_{W_0}\le I_0$.
Since $\nrm{\partial_\theta q}_{W_0}\le RP'$,
differentiating $d_\partial$ gives the bound $D_\theta$.
Differentiating $d_\partial d_\partial^{-1}=1$ then gives
$\nrm{d_\partial^{-1}}_{W_1}\le I_1$.
This is the weaker-norm reciprocal argument in
Proposition~\ref{r2:prop:coupled}; smallness in $W_1$ is unnecessary.

By~\eqref{gd:eq:Dirichlet-recipe} and~\eqref{gd:eq:JW},
$\nrm{\mathfrak a}_{W_1}\le I_1N_W\nrm f_Y$.
Equation~\eqref{gd:eq:coupled-recipe} therefore gives
\begin{equation}\label{gd:eq:A}
 B_g=J_W+M_TI_1N_W,\qquad
 B_p=(P+RP')I_1N_W,\qquad A_*=B_g+\sigma B_p,\qquad
 \nrm A_{Y\to X}\le A_* .
\end{equation}
Here $\delta p=(wp_0\mathfrak a)'$ has the norm of $p_0\mathfrak a$ in $W_1$.
For an arbitrary input variation, $v=\eta/p_0$ gives
\begin{equation}\label{gd:eq:Jv}
 \nrm{v'}_W\le J_v\nrm{\delta p}_W,\qquad
 J_v=q_0^{-1}+RP'q_0^{-2},
\end{equation}
using $\nrm\eta_W\le R\nrm{\delta p}_W$.
The material correction has the bound
\begin{equation}\label{gd:eq:epsilon}
 \nrm{\Ecal}_{X\to Y}\le\varepsilon_{\mathcal E}
 =\frac{\kappa}{\sigma}
       \left[(RC_x^{\rm pol}+C_y^{\rm pol})E_d+3E_0\right]J_v .
\end{equation}
The factor three is $2+1$ from the last two terms of
\eqref{gd:eq:correction}. The transport terms use
$\im v\,E_y=(\im v/y)yE_y$.
Since $A\widetilde D=I_X$, the left Newton defect is
\begin{equation}\label{gd:eq:Newton}
 \nrm{AF(x_0)}_X\le Y_0=A_*E_0,\qquad
 \nrm{I_X-ADF(x_0)}_{X\to X}\le Z=A_*\varepsilon_{\mathcal E} .
\end{equation}
The field right sides need only be continuous.
Their Dirichlet potentials have $C^1$ traces by the preceding
Poisson and elliptic bounds, whereas $T$ is analytic near the closed
ball. Thus the identities apply to every element of $X$ and $Y$.

\begin{table}[htbp]
\centering\small
\def\gdRow#1#2#3#4#5#6#7#8{$#1$ & #5\\}
\begin{tabular}{@{}crrrrr@{}}
\toprule
$n$ & $A_*$ & $Y_0$ & $Z$ & $4C_2Y_0$ & $\mathfrak c(r_*)$\\
\midrule
\gdData
\bottomrule
\end{tabular}
\caption{Corrected inverse and contraction bounds.
$\mathfrak c$ is defined in~\eqref{gd:eq:radii}.
The quantity $4C_2Y_0$ is defined by~\eqref{gd:eq:Newton}
and~\eqref{gd:eq:radii}.}
\label{gd:tab:Newton}
\end{table}

\section{Quartic bounds and the exact zero}\label{gd:sec:existence}

For a variation $(h,k)\in X$ write $u_h=K_nh$ and
$q_k=\im\int_0^wk(z)\,dz/y$.
Lemma~\ref{gd:lem:connection} and the axis multiplier estimates give
\begin{equation}\label{gd:eq:shape-increments}
 \begin{gathered}
 \nrm{q_k}_Y\le\kappa\nrm k_W,\quad
 \nrm{(q_k)_{x_1}}_{\mathcal Y_n}\le\kappa c_a\nrm k_W,\quad
 \nrm{(q_k)_y/y}_Y\le\kappa c_v\nrm k_W,\\
 \nrm{p_0\bar k+\bar p_0k}_Y\le2\kappa P\nrm k_W,\qquad
 \nrm{|k|^2}_Y\le\kappa\nrm k_W^2 .
 \end{gathered}
\end{equation}
The last two estimates multiply in the raw angular algebra before
conversion. These bounded linear operators, together with
Lemma~\ref{gd:lem:Poisson}, make $F:X\to Y$ a continuous quartic
polynomial on this fixed Banach space.

To bound its remainder, put $P_k=p_0\bar k+\bar p_0k$.
The terms of degree two, after removal of $E+DF(x_0)(h,k)$,
are
\[
 \begin{split}
 b\mathcal N_2={}&q_kh+
 a\{(q_k)_{x_1}(u_h)_{x_1}+((q_k)_y/y)y(u_h)_y\}\\
 &+\mu\{(qP_k+q_k|p_0|^2)u_h+(q|k|^2+q_kP_k)U_0\}.
 \end{split}
\]
The terms of degrees three and four are
\[
 b\mathcal N_3=\mu\{(q|k|^2+q_kP_k)u_h+q_k|k|^2U_0\},
 \qquad b\mathcal N_4=\mu q_k|k|^2u_h .
\]
Projection to the field has norm at most one and projection to the
shape has norm at most $1/\sigma$ in $X$.
The algebra, Poisson, and centre estimates therefore give
\begin{equation}\label{gd:eq:c234}
 \begin{split}
 \mathsf k_2={}&\frac1b\biggl[
 \frac{\kappa}{\sigma}
 +\frac{a\kappa(c_aC_x^K+c_vC_y^K)}{\sigma}
 +\frac{\mu(2\kappa Q_*P+\kappa P_*)C_K}{\sigma}\\
 &\hspace{24mm}
       +\frac{\mu(\kappa Q_*+2\kappa^2P)U_c}{\sigma^2}\biggr],\\
 \mathsf k_3={}&\frac\mu b\left[
       \frac{(\kappa Q_*+2\kappa^2P)C_K}{\sigma^2}
       +\frac{\kappa^2U_c}{\sigma^3}\right],\\
 \mathsf k_4={}&\frac{\mu\kappa^2C_K}{b\sigma^3}.
 \end{split}
\end{equation}
In particular $\nrm{\mathcal N_j(h,k)}_Y\le
\mathsf k_j\nrm{(h,k)}_X^j$.
Each displayed product also defines a multilinear form with this same
bound on the product of input norms. Symmetrisation averages over the
positions of the field and shape inputs, and introduces no larger
constant. Differentiating the diagonal form consequently bounds
its derivative by $j\mathsf k_js^{j-1}$ on $\nrm{(h,k)}_X\le s$,
as in the polarisation argument of Section~\ref{r2:sec:existence}.

Set $C_j=A_*\mathsf k_j$ and define
\begin{equation}\label{gd:eq:radii}
 \begin{split}
 \mathfrak p(s)&=Y_0+(Z-1)s+C_2s^2+C_3s^3+C_4s^4,\\
 \mathfrak c(s)&=Z+2C_2s+3C_3s^2+4C_4s^3 .
 \end{split}
\end{equation}

\begin{proposition}\label{gd:prop:zero}
If $\mathfrak p(r_*)<0$ and $\mathfrak c(r_*)<1$, then
$F$ has exactly one zero in
$\overline B_X(x_0,r_*)$.
\end{proposition}
\begin{proof}
The proof of Lemma~\ref{r2:lem:contraction} applies to the additional
quartic term. For $\mathcal T(x)=x-AF(x)$, the estimates give
\[
 \nrm{\mathcal T(x_0+h)-x_0}_X
 \le Y_0+Z\nrm h_X+\sum_{j=2}^4C_j\nrm h_X^j .
\]
The first inequality makes the closed ball invariant.
Its derivative is bounded there by $\mathfrak c(r_*)<1$,
so completeness gives a unique fixed point.
The two-sided inverse $A$ is injective, hence $AF(x)=0$ implies
$F(x)=0$. Every zero is a fixed point.
\end{proof}

\begin{table}[htbp]
\centering\small
\def\gdRow#1#2#3#4#5#6#7#8{$#1$ & #6\\}
\begin{tabular}{@{}crrrrr@{}}
\toprule
$n$ & $E_0$ & $C_2$ & $C_3$ & $C_4$ & $\mathfrak p(r_*)$\\
\midrule
\gdData
\bottomrule
\end{tabular}
\caption{The full residual, homogeneous remainder constants and
negative radii polynomial, at the exact radius of
Table~\ref{gd:tab:inputs}.}\label{gd:tab:nonlinear}
\end{table}

Tables~\ref{gd:tab:Newton} and~\ref{gd:tab:nonlinear} give
$\mathfrak p(r_*)<0$ and $\mathfrak c(r_*)<10^{-5}$,
so Proposition~\ref{gd:prop:zero} applies.

\section{Geometry on the whole existence ball}\label{gd:sec:geometry}

Write $\psi_0=\sum_{j\ {\rm odd}}c_jw^j$ and let
$\psi=\psi_0+\delta\psi$ be any shape in the existence ball.
Then $\nrm{\delta p}_W\le s_*=r_*/\sigma$.
Define the unweighted centre sums
\begin{equation}\label{gd:eq:geometry-sums}
 \begin{gathered}
 S_1=\sum_{j>1}j|c_j|,\qquad
 S_2=\sum_{j>1}j(j-1)|c_j|,\\
 S_T=\sum_{j>1}j^2|c_j|=S_1+S_2,\qquad
 H_R=\frac{R}{(R-1)^2},\\
 \beta_0=2b-P-s_*,\quad L_*=b-S_1-s_*,\quad
 \chi_*=b-S_1-S_2-(1+H_R)s_*,\\
 \vartheta_3=|c_3|-\frac{s_*}{3R^2}.
 \end{gathered}
\end{equation}
Table~\ref{gd:tab:geometry} gives strict positive lower bounds.
Since $P-b\ge S_1$, $\beta_0>0$ also implies $L_*>0$.
The two sums $S_T$ and $S_1+S_2$ are identical, so the meridian and
transverse estimates below use one common numerator $\chi_*$.

\begin{table}[htbp]
\centering
\def\gdRow#1#2#3#4#5#6#7#8{$#1$ & #7\\}
\begin{tabular}{@{}crrr@{}}
\toprule
$n$ & $\beta_0$ & $\chi_*$ & $\vartheta_3$\\
\midrule
\gdData
\bottomrule
\end{tabular}
\caption{Whole-ball geometry: analytic univalence, the common
curvature numerator and the absolute $w^3$ coefficient. The numerator
$\chi_*$ is the common numerator of
\eqref{gd:eq:meridian-curvature} and~\eqref{gd:eq:transverse-numerator}.}
\label{gd:tab:geometry}
\end{table}

\subsection{Univalence and analytic lift}\label{gd:sec:univalence}
On $|w|\le R$,
\[
 \re p(w)\ge b-(P-b)-s_*=\beta_0>0.
\]
On the closed unit disc the sharper coefficient estimate is
\begin{equation}\label{gd:eq:unit-disc-tangent}
 \re\psi'(w)\ge b-S_1-s_*=L_*>0\qquad(|w|\le1).
\end{equation}
The line-segment proof in Section~\ref{r2:sec:existence} applies verbatim:
for distinct points $w_1,w_2$ of the disc $|w|<R$,
\[
 \frac{\psi(w_2)-\psi(w_1)}{w_2-w_1}
 =\int_0^1p(w_1+t(w_2-w_1))\,dt
\]
has positive real part. Thus $\psi$ is univalent on $|w|<R$.
Also $q=\int_0^1\re p(x_1+isy)\,ds>0$ there, so $\psi$ preserves
the upper half-disc. The real-coefficient identities express
$\re\psi$ and $q$ as convergent power series in $x_1,y^2$.
Replacing $y^2$ by $|x'|^2$ makes $\Theta_\psi$ analytic across the axis.

Off the axis the meridian differential and the $n-2$ transverse
differentials give $\det D\Theta_\psi=|p|^2q^{n-2}>0$.
On the axis its differential is
$\diag(p(x_1),p(x_1)I_{n-1})$, also nonsingular.
If two lifted points have the same image, their meridian coordinates
have the same image under $\psi$; univalence gives equality of $x_1,y$,
and positivity of $q$ then gives equality of $x'$.
Thus the lift is an analytic diffeomorphism on a neighbourhood of
$\overline{B^n}$. Its restriction to the sphere gives a real-analytic
boundary diffeomorphic to $S^{n-1}$.
Oddness and reality give the required $O(n-1)\times\Z_2$ symmetry.

\subsection{Meridian and transverse curvatures}\label{gd:sec:curvatures}
For the perturbation,
\[
 \sup_{|w|\le1}|\delta p'|
 \le\sup_{m\ge1}mR^{-m}\nrm{\delta p}_W
 \le H_Rs_* .
\]
Therefore $|p|\ge L_*$ and $|wp'|\le S_2+H_Rs_*$ on the closed disc.
The curvature formula~\eqref{r2:eq:curvature-formula} gives
\begin{equation}\label{gd:eq:meridian-curvature}
 \re(1+wp'/p)\ge\frac{\chi_*}{L_*}>0,\qquad
 \kappa_{\rm mer}=\frac{\re(1+wp'/p)}{|p|}>0 .
\end{equation}
The simple meridian boundary consequently encloses a strictly convex
planar domain symmetric in both axes.

For the transverse curvatures and the graph property, write
$\delta\psi=\sum_{j\ {\rm odd}}\delta c_jw^j$ and parameterise the upper arc as
$\psi(e^{i\theta})=(X(\theta),Y_\psi(\theta))$, $0<\theta<\pi$.
Its outward normal is $(Y_\psi',-X')/|p|$ and speed is $|p|$.
Using $U_m(\cos\theta)=\sin((m+1)\theta)/\sin\theta$ gives
\begin{equation}\label{gd:eq:transverse-numerator}
 -\frac{X'}{\sin\theta}
 =\sum_{j\ {\rm odd}}j(c_j+\delta c_j)U_{j-1}(\cos\theta)
 \ge b-S_T-(1+H_R)s_*=\chi_* .
\end{equation}
Indeed $|U_{j-1}|\le j$, and
\[
 \sum_{j\ {\rm odd}}j^2|\delta c_j|
 =\sum_{m\ {\rm even}}(m+1)|\delta p_m|
 \le(1+H_R)\nrm{\delta p}_W.
\]
Since $Y_\psi/\sin\theta=q>0$, each of the $n-2$ transverse principal
curvatures of the unscaled body is
\begin{equation}\label{gd:eq:transverse-curvature}
 \kappa_\perp=\frac{-X'}{|p|Y_\psi}
 =\frac{-X'/\sin\theta}{|p|(Y_\psi/\sin\theta)}
 \ge\frac{\chi_*}{(b+S_1+s_*)^2}>0.
\end{equation}
Both denominators have the displayed upper bound by the coefficient
and integral estimates for $p,q$.
The quotients extend to the poles analytically.
At $\theta=0$, $Y_\psi'=p(1)>0$ and
$-X'/\sin\theta\to-X''(0)=p(1)+p'(1)$, so
$\kappa_\perp$ tends to
$(p(1)+p'(1))/p(1)^2$, the meridian curvature there.
The other pole has the same conclusion by symmetry.
Thus every principal curvature is positive on the whole boundary.
Scaling by $\rho$ divides all principal curvatures by $\rho$.

\subsection{Global convexity and the non-ball coefficient}
\label{gd:sec:global}
By~\eqref{gd:eq:transverse-numerator}, $X'(\theta)<0$ on the open
upper arc. It is therefore a graph $Y_\psi=f(X)$ between its two poles.
The curvature numerator of this oriented graph is
$X'Y_\psi''-Y_\psi'X''=(X')^3f''$; positive meridian curvature and $X'<0$
give $f''<0$ in the interior. Continuity extends its concavity to the
closed interval, where $f$ vanishes at the endpoints and is positive
inside. The lifted domain is exactly
\[
 \Theta_\psi(B^n)=\{(X_1,X_\perp):X_1\text{ between the poles},\
                 X_\perp\in\R^{n-1},\ |X_\perp|<f(X_1)\}.
\]
For two interior points $\mathbf X=(X_1,X_\perp)$ and
$\tilde{\mathbf X}=(\tilde X_1,\tilde X_\perp)$ and $0<t<1$,
the triangle inequality and concavity give
\[
 \begin{split}
 |(1-t)X_\perp+t\tilde X_\perp|
 &\le(1-t)|X_\perp|+t|\tilde X_\perp|\\
 &<(1-t)f(X_1)+tf(\tilde X_1)\\
 &\le f((1-t)X_1+t\tilde X_1).
 \end{split}
\]
Hence the body is convex globally.
Its smooth boundary has positive principal curvatures by
\eqref{gd:eq:meridian-curvature} and~\eqref{gd:eq:transverse-curvature},
so it contains no boundary line segment and is strictly convex.
This is the solid-of-revolution argument used in
Sections~\ref{r5:sec-zero} and~\ref{r7:sec-zero}.

Finally the coefficient estimate gives
$|c_3+\delta c_3|\ge\vartheta_3>0$.
A ball invariant under $O(n-1)$ and first-coordinate reflection has
centre zero. Its meridian section would be a centred disc.
The odd conformal map fixing zero with positive derivative there
would then be linear, by uniqueness of the normalised conformal map,
contradicting its nonzero $w^3$ coefficient.
This excludes balls throughout the shape projection of the existence
ball. The origin $\Theta_\psi(0)$ is an interior point of the convex
domain. Every tangent hyperplane supports the domain and therefore
does not contain the origin. Every ray from the origin meets the
boundary exactly once and transversally, so $\Omega$ is strictly
star-shaped with respect to the origin.

\section{The Schiffer field and the Pompeiu conclusion}
\label{gd:sec:conclusion}

\begin{proof}[Proof of Theorem~\ref{gd:thm:main}]
For each $n\in\gdDimensions$, the centre satisfies
$c=\rho b>20.69$, $q_*>0.9548$ and $\gamma>299.4$.
The window of Lemma~\ref{gd:lem:coverage} has the state count of
Table~\ref{gd:tab:inverse}. The hypotheses~\eqref{gd:eq:Schur-gates} hold with
$\epsilon_C<2.7\cdot10^{-10}$, $\epsilon_{\rm S}+\eta_{\rm S}<0.0242$
and $B_{\rm pol}\delta_V<1.1\cdot10^{-4}$, so
Proposition~\ref{gd:prop:Schur} gives $B_{L^2}$.
With the tabulated $L_0$, one has $\beta<0.4987$ and
$q_b<0.3412<b$. Thus~\eqref{gd:eq:BY} and
Proposition~\ref{gd:prop:original-inverse} give $B_Y,J_W,N_W$.
For both strip radii, $d_s<0.13b$ in~\eqref{gd:eq:strip-parameters};
also $q_0>0.9478$ and $d_*>387.0$ in~\eqref{gd:eq:boundary-reciprocal}.
Proposition~\ref{gd:prop:coupled} therefore applies with the bound
$A_*$ of Table~\ref{gd:tab:Newton}. The strict radii inequalities
in Tables~\ref{gd:tab:Newton} and~\ref{gd:tab:nonlinear} give the unique
zero $(g,p)$ in $\overline B_X(x_0,r_*)$ by
Proposition~\ref{gd:prop:zero}.
Put $U=1+K_ng$. Lemma~\ref{gd:lem:Poisson} gives $U\in W^{2,p}(B^n)$
for every finite $p$, with $U=1$ and $\partial_rU=0$ on the sphere.
Its tangential derivatives also vanish there, so $\nabla U=0$.
Since $U\in W^{2,p}$, $F(g,p)=0$ says $BU=0$ almost everywhere.
Multiplication by $bq^{a-1}$ and integration by parts give
\[
 \int_{B^n}q^a(\nabla U\cdot\nabla\varphi-\mu|p|^2U\varphi)\,dx=0
 \qquad(\varphi\in C_c^\infty(B^n)).
\]
The analytic diffeomorphism of Section~\ref{gd:sec:univalence} and
the change of variables of Lemma~\ref{r5:lem-lift}(b), with $n-2$
transverse directions, give
$u_0=U\circ\Theta_\psi^{-1}\in C^1(\overline{\Theta_\psi(B^n)})$,
a weak solution of $(\Delta+\mu)u_0=0$ with $u_0=1$ and
$\nabla u_0=0$ on the boundary. Set
\[
 \Omega=\rho\Theta_\psi(B^n),\qquad u(X)=u_0(X/\rho).
\]
Since $\mu=\rho^2$, $u\in C^1(\overline\Omega)$ solves
$\Delta u+u=0$ weakly with $u=1$, $\nabla u=0$ on the real-analytic
boundary. Lemma~\ref{r3:lem-regularity} gives a real-analytic extension
to a neighbourhood of $\overline\Omega$, solving~\eqref{gd:eq:schiffer}.
It cannot be constant, because its boundary value would make it the
constant one, which does not solve $(\Delta+1)u=0$.
Sections~\ref{gd:sec:univalence}--\ref{gd:sec:global} give all geometric
assertions and the symmetry.

The Green identity of Section~\ref{sec:pompeiu}
applies in every dimension. For a unit vector $\mathbf k$ put
$v(X)=e^{i\mathbf k\cdot X}$. Then
\[
 0=\int_\Omega(u\Delta v-v\Delta u)
   =\int_{\partial\Omega}\partial_{\mathbf n} v
   =-\int_\Omega e^{i\mathbf k\cdot X}\,dX .
\]
Thus $\widehat{\one_\Omega}$ vanishes on the unit sphere.
For a rigid motion $X\mapsto QX+A$,
\[
 \int_{Q\Omega+A}e^{iX_1}\,dX
   =e^{iA_1}\int_\Omega e^{i(Q^Te_1)\cdot X}\,dX=0 .
\]
Taking real parts gives zero integrals of the nonzero continuous
function $\cos X_1$ over every rigid motion of $\Omega$.
Hence $\Omega$ fails the Pompeiu property.
\end{proof}
